\Article{1.}
Investigations, in which the directions of various straight lines in space are to be
considered, attain a high degree of clearness and simplicity if we employ, as an auxiliary,
a sphere of unit radius described about an arbitrary centre, and suppose the
different points of the sphere to represent the directions of straight lines parallel to
the radii ending at these points. As the position of every point in space is determined
by three coordinates, that is to say, the distances of the point from three mutually
perpendicular fixed planes, it is necessary to consider, first of all, the directions of the
axes perpendicular to these planes. The points on the sphere, which represent these
directions, we shall denote by $(1)$,~$(2)$,~$(3)$. The distance of any one of these points
from either of the other two will be a quadrant; and we shall suppose that the directions
of the axes are those in which the corresponding coordinates increase.


\Article{2.}
It will be advantageous to bring together here some propositions which are frequently
used in questions of this kind.

\Par{I.} The angle between two intersecting straight lines is measured by the arc
between the points on the sphere which correspond to the directions of the lines.

\Par{II.}\Note{1} The orientation of any plane whatever can be represented by the great circle
on the sphere, the plane of which is parallel to the given plane.
\PageSep{4}

\Par{III.} The angle between two planes is equal to the spherical angle between the
great circles representing them, and, consequently, is also measured by the arc intercepted
between the poles of these great circles. And, in like manner, the angle of inclination
of a straight line to a plane is measured by the arc drawn from the point which
corresponds to the direction of the line, perpendicular to the great circle which represents
the orientation of the plane.

\Par{IV.} Letting $x$,~$y$,~$z$; $x'$,~$y'$,~$z'$ denote the coordinates of two points, $r$~the distance
between them, and $L$~the point on the sphere which represents the direction of the line
drawn from the first point to the second, we shall have
\begin{align*}
x' &= x + r \cos(1)L\Add{,} \\
y' &= y + r \cos(2)L\Add{,} \\
z' &= z + r \cos(3)L\Add{.}
\end{align*}

\Par{V.} From this it follows at once that, generally,\Note{2}
\[
\cos^{2}(1)L + \cos^{2}(2)L + \cos^{2}(3)L = 1\Add{,}
\]
and also, if $L'$~denote any other point on the sphere,
\[
\cos(1)L·\cos(1)L' + \cos(2)L·\cos(2)L' + \cos(3)L·\cos(3)L'
  = \cos LL'.
\]

\Par{VI.}\Note{4} \begin{Theorem}
If $L$, $L'$, $L''$, $L'''$ denote four points on the sphere, and $A$~the angle
which the arcs $LL'$, $L''L'''$ make at their point of intersection, then we shall have
\[
\cos LL''·\cos L'L''' - \cos LL'''·\cos L'L''
  = \sin LL'·\sin L''L'''·\cos A\Add{.}
\]
\end{Theorem}

\textit{Demonstration.} Let $A$ denote also the point of intersection itself, and set
\[
AL = t,\quad AL' = t',\quad AL'' = t'',\quad AL''' = t'''\Add{.}
\]
Then we shall have
\begin{alignat*}{7}
&\cos L L''  &&= \cos t\Chg{·}{}&& \cos t''  &&+ \sin t  && \sin t''  && \cos A\Add{,} \\
&\cos L'L''' &&= \cos t' && \cos t''' &&+ \sin t' && \sin t''' && \cos A\Add{,} \\
&\cos L L''' &&= \cos t  && \cos t''' &&+ \sin t  && \sin t''' && \cos A\Add{,} \\
&\cos L'L''  &&= \cos t' && \cos t''  &&+ \sin t' && \sin t''  && \cos A\Add{;}
\end{alignat*}
and consequently,\Note{3}
\begin{multline*}
\cos LL''·\cos L'L''' - \cos LL'''·\cos L'L'' \\
\begin{aligned}
&= \cos A (\cos t  \cos t''  \sin t' \sin t'''
         + \cos t' \cos t''' \sin t  \sin t'' \\
&\qquad  - \cos t  \cos t''' \sin t' \sin t''
         - \cos t' \cos t''  \sin t  \sin t''') \\
&= \cos A (\cos t   \sin t'   - \sin t   \cos t')
          (\cos t'' \sin t''' - \sin t'' \cos t''')\NoteMark \\
&= \cos A·\sin (t' - t)·\sin (t''' - t'') \\
&= \cos A·\sin LL'·\sin L''L'''\Add{.}
\end{aligned}
\end{multline*}
\PageSep{5}

But as there are for each great circle two branches going out from the point~$A$,
these two branches form at this point two angles whose sum is~$180°$. But our analysis
shows that those branches are to be taken whose directions are in the sense from the
point $L$~to~$L'$, and from the point $L''$~to~$L'''$; and since great circles intersect in two
points, it is clear that either of the two points can be chosen arbitrarily. Also, instead
of the angle~$A$, we can take the arc between the poles of the great circles of which the
arcs $LL'$,~$L''L'''$ are parts. But it is evident that those poles are to be chosen which
are similarly placed with respect to these arcs; that is to say, when we go from $L$~to~$L'$
and from $L''$~to~$L'''$, both of the two poles are to be on the right, or both on the left.

\Par{VII.} Let $L$,~$L'$,~$L''$ be the three points on the sphere and set, for brevity,
\begin{alignat*}{6}
&\cos (1)L   &&= x,\quad&& \cos (2)L   &&= y,\quad&& \cos (3)L   &&= z\Add{,} \\
&\cos (1)L'  &&= x',    && \cos (2)L'  &&= y',    && \cos (3)L'  &&= z'\Add{,} \\
&\cos (1)L'' &&= x'',   && \cos (2)L'' &&= y'',   && \cos (3)L'' &&= z''\Add{;} \\
\end{alignat*}
and also
\[
x y' z'' + x' y'' z + x'' y z' - x y'' z' - x' y z'' - x'' y' z = \Delta\Add{.}
\]
Let $\lambda$~denote the pole of the great circle of which $LL'$~is a part, this pole being the one
that is placed in the same position with respect to this arc as the point~$(1)$ is with
respect to the arc~$(2)(3)$. Then we shall have, by the preceding theorem,
\[
y z' - y' z = \cos (1)\lambda·\sin (2)(3)·\sin LL',
\]
or, because $(2)(3) = 90°$,
\begin{align*}
y z' - y' z &= \cos (1)\lambda·\sin LL', \\
\intertext{and similarly,}
z x' - z' x &= \cos (2)\lambda·\sin LL'\Add{,} \\
x y' - x' y &= \cos (3)\lambda·\sin LL'\Add{.}
\end{align*}
Multiplying these equations by $x''$,~$y''$,~$z''$ respectively, and adding, we obtain, by means
of the second of the theorems deduced in~V,
\[
\Delta = \cos \lambda L''·\sin LL'\Add{.}
\]
Now there are three cases to be distinguished. \emph{First}, when $L''$~lies on the great circle
of which the arc~$LL'$ is a part, we shall have $\lambda L'' = 90°$, and consequently, $\Delta = 0$.
If $L''$~does not lie on that great circle, the \emph{second} case will be when $L''$~is on the same
side as~$\lambda$; the \emph{third} case when they are on opposite sides. In the last two cases the
points $L$,~$L'$,~$L''$ will form a spherical triangle, and in the second case these points will lie
in the same order as the points $(1)$,~$(2)$,~$(3)$, and in the opposite order in the third case.
\PageSep{6}
Denoting the angles of this triangle simply by $L$,~$L'$,~$L''$ and the perpendicular drawn on
the sphere from the point~$L''$ to the side~$LL'$ by~$p$, we shall have
\[
\sin p = \sin L·\sin LL'' = \sin L'·\sin L' L'',
\]
and
\[
\lambda L'' = 90° \mp p,
\]
the upper sign being taken for the second case, the lower for the third. From this
it follows that
\begin{align*}
±\Delta &= \sin L·\sin LL'·\sin LL''
   = \sin L'·\sin LL'·\sin L'L'' \\
  &= \sin L''·\sin LL''·\sin L'L''\Add{.}
\end{align*}
Moreover, it is evident that the first case can be regarded as contained in the second or
third, and it is easily seen that the expression~$±\Delta$ represents six times the volume of
the pyramid formed by the points $L$,~$L'$,~$L''$ and the centre of the sphere. Whence,
finally, it is clear that the expression~$±\frac{1}{6}\Delta$ expresses generally the volume of any
pyramid contained between the origin of coordinates and the three points whose coordinates
are $\Typo{z}{x}$,~$y$,~$z$; $x'$,~$y'$,~$z'$; $x''$,~$y''$,~$z''$.\Note{5}


\Article{3.}

A curved surface is said to possess continuous curvature at one of its points~$A$, if the
directions of all the straight lines drawn from $A$ to points of the surface at an infinitely
small distance from~$A$ are deflected infinitely little from one and the same plane passing
through~$A$. This plane is said to \emph{touch} the surface at the point~$A$. If this condition is
not satisfied for any point, the continuity of the curvature is here interrupted, as happens,
for example, at the vertex of a cone. The following investigations will be restricted to
such surfaces, or to such parts of surfaces, as have the continuity of their curvature
nowhere interrupted. We shall only observe now that the methods used to determine
the position of the tangent plane lose their meaning at singular points, in which the
continuity of the curvature is interrupted, and must lead to indeterminate solutions.


\Article{4.}

The orientation of the tangent plane is most conveniently studied by means of the
direction of the straight line normal to the plane at the point~$A$, which is also called the
normal to the curved surface at the point~$A$. We shall represent the direction of this
normal by the point~$L$ on the auxiliary sphere, and we shall set
\[
\cos (1)L = X,\quad \cos (2)L = Y,\quad \cos (3)L = Z;
\]
and denote the coordinates of the point~$A$ by $x$,~$y$,~$z$. Also let $x + dx$, $y + dy$, $z + dz$
be the coordinates of another point~$A'$ on the curved surface; $ds$~its distance from~$A$,
\PageSep{7}
which is infinitely small; and finally, let $\lambda$ be the point on the sphere representing the
direction of the element~$AA'$. Then we shall have
\[
dx = ds·\cos (1)\lambda,\quad
dy = ds·\cos (2)\lambda,\quad
dz = ds·\cos (3)\lambda
\]
and, since $\lambda L$~must be equal to~$90°$,
\[
X\cos (1)\lambda + Y\cos (2)\lambda + Z\cos (3)\lambda = 0\Add{.}
\]
By combining these equations we obtain
\[
X\, dx + Y\, dy + Z\, dz = 0.
\]

There are two general methods for defining the nature of a curved surface. The
\emph{first} uses the equation between the coordinates $x$,~$y$,~$z$, which we may suppose reduced to
the form $W = 0$, where $W$~will be a function of the indeterminates $x$,~$y$,~$z$. Let the complete
differential of the function~$W$ be
\[
dW = P\, dx + Q\, dy + R\, dz
\]
and on the curved surface we shall have
\[
P\, dx + Q\, dy + R\, dz = 0\Add{,}
\]
and consequently,
\[
P \cos (1)\lambda + Q \cos (2)\lambda + R \cos (3)\lambda = 0\Add{.}
\]
Since this equation, as well as the one we have established above, must be true for the
directions of all elements~$ds$ on the curved surface, we easily see that $X$,~$Y$,~$Z$ must be
proportional to $P$,~$Q$,~$R$ respectively, and consequently, since\Note{6}
\[
X^{2} + Y^{2} + Z^{2} = 1,\NoteMark
\]
we shall have either
\begin{align*}
X &= \frac{P}{\Sqrt{P^{2} + Q^{2} + R^{2}}}, &
Y &= \frac{Q}{\Sqrt{P^{2} + Q^{2} + R^{2}}}, &
Z &= \frac{R}{\Sqrt{P^{2} + Q^{2} + R^{2}}} \\
\intertext{or}
X &= \frac{-P}{\Sqrt{P^{2} + Q^{2} + R^{2}}}, &
Y &= \frac{-Q}{\Sqrt{P^{2} + Q^{2} + R^{2}}}, &
Z &= \frac{-R}{\Sqrt{P^{2} + Q^{2} + R^{2}}}\Add{.}
\end{align*}

The \emph{second}\Note{7} method expresses the coordinates in the form of functions of two variables,
$p$,~$q$. Suppose that differentiation of these functions gives
\begin{alignat*}{2}
dx &= a\, dp &&+ a'\, dq\Add{,} \\
dy &= b\, dp &&+ b'\, dq\Add{,} \\
dz &= c\, dp &&+ c'\, dq\Add{.}
\end{alignat*}
\PageSep{8}
Substituting these values in the formula given above, we obtain
\[
(aX + bY + cZ)\, dp + (a'X + b'Y + c'Z)\, dq = 0\Add{.}
\]
Since this equation must hold independently of the values of the differentials $dp$,~$dq$,
we evidently shall have
\[
aX + bY + cZ = 0,\quad a'X + b'Y + c'Z = 0\Add{.}
\]
From this we see that $X$,~$Y$,~$Z$ will be proportioned to the quantities
\[
bc' - cb',\quad ca' - ac',\quad ab' - ba'\Add{.}
\]
Hence, on setting, for brevity,
\[
\Sqrt{(bc' - cb')^{2} + (ca' - ac')^{2} + (ab' - ba')^{2}} = \Delta\Add{,}
\]
we shall have either
\begin{align*}
X &= \frac{bc' - cb'}{\Delta},\quad&
Y &= \frac{ca' - ac'}{\Delta},\quad&
Z &= \frac{ab' - ba'}{\Delta}
\intertext{or}
X &= \frac{cb' - bc'}{\Delta},\quad&
Y &= \frac{ac' - ca'}{\Delta},\quad&
Z &= \frac{ba' - ab'}{\Delta}\Add{.}
\end{align*}

With these two general methods is associated a \emph{third}, in which one of the coordinates,
$z$,~say, is expressed in the form of a function of the other two, $x$,~$y$. This method is
evidently only a particular case either of the first method, or of the second. If we set
\[
dz = t\, dx + u\, dy
\]
we shall have either
\begin{align*}
X &= \frac{-t}{\Sqrt{1 + t^{2} + u^{2}}}, &
Y &= \frac{-u}{\Sqrt{1 + t^{2} + u^{2}}}, &
Z &= \frac{ 1}{\Sqrt{1 + t^{2} + u^{2}}} \\
\intertext{or}
X &= \frac{ t}{\Sqrt{1 + t^{2} + u^{2}}}, &
Y &= \frac{ u}{\Sqrt{1 + t^{2} + u^{2}}}, &
Z &= \frac{-1}{\Sqrt{1 + t^{2} + u^{2}}}\Add{.}
\end{align*}


\Article{5.}

The two solutions found in the preceding article evidently refer to opposite points of
the sphere, or to opposite directions, as one would expect, since the normal may be drawn
toward either of the two sides of the curved surface. If we wish to distinguish between
the two regions bordering upon the surface, and call one the exterior region and the other
the interior region, we can then assign to each of the two normals its appropriate solution
by aid of the theorem derived in \Art{2}~(VII), and at the same time establish a criterion
for distinguishing the one region from the other.
\PageSep{9}

In the first method, such a criterion is to be drawn from the sign of the quantity~$W$.
Indeed, generally speaking, the curved surface divides those regions of space in which $W$
keeps a positive value from those in which the value of~$W$ becomes negative. In fact, it
is easily seen from this theorem that, if $W$ takes a positive value toward the exterior
region, and if the normal is supposed to be drawn outwardly, the first solution is to be
taken. Moreover, it will be easy to decide in any case whether the same rule for the
sign of~$W$ is to hold throughout the entire surface, or whether for different parts there
will be different rules. As long as the coefficients $P$,~$Q$,~$R$ have finite values and do not
all vanish at the same time, the law of continuity will prevent any change.

If we follow the second method, we can imagine two systems of curved lines on the
curved surface, one system for which $p$~is variable, $q$~constant; the other for which $q$~is
variable, $p$~constant. The respective positions of these lines with reference to the exterior
region will decide which of the two solutions must be taken. In fact, whenever
the three lines, namely, the branch of the line of the former system going out from the
point~$A$ as $p$~increases, the branch of the line of the latter system going out from the point
$A$ as $q$~increases, and the normal drawn toward the exterior region, are \emph{similarly} placed as
the $x$,~$y$,~$z$ axes respectively from the origin of abscissas (\eg, if, both for the former
three lines and for the latter three, we can conceive the first directed to the left, the
second to the right, and the third upward), the first solution is to be taken. But whenever
the relative position of the three lines is opposite to the relative position of the
$x$,~$y$,~$z$ axes, the second solution will hold.

In the third method, it is to be seen whether, when $z$~receives a positive increment, $x$~and~$y$
remaining constant, the point crosses toward the exterior or the interior region.
In the former case, for the normal drawn outward, the first solution holds; in the latter
case, the second.


\Article{6.}

Just as each definite point on the curved surface is made to correspond to a definite
point on the sphere, by the direction of the normal to the curved surface which is transferred
to the surface of the sphere, so also any line whatever, or any figure whatever, on
the latter will be represented by a corresponding line or figure on the former. In the
comparison of two figures corresponding to one another in this way, one of which will be
as the map of the other, two important points are to be considered, one when quantity
alone is considered, the other when, disregarding quantitative relations, position alone
is considered.

The first of these important points will be the basis of some ideas which it seems
judicious to introduce into the theory of curved surfaces. Thus, to each part of a curved
\PageSep{10}
surface inclosed within definite limits we assign a \emph{total} or \emph{integral curvature}, which is
represented by the area of the figure on the sphere corresponding to it. From this
integral curvature must be distinguished the somewhat more specific curvature which we
shall call the\Note{8} \emph{measure of curvature}. The latter refers to a \emph{point} of the surface, and shall
denote the quotient obtained when the integral curvature of the surface element about
a point is divided by the area of the element itself; and hence it denotes the ratio of the
infinitely small areas which correspond to one another on the curved surface and on the
sphere. The use of these innovations will be abundantly justified, as we hope, by what
we shall explain below. As for the terminology, we have thought it especially desirable
that all ambiguity be avoided. For this reason we have not thought it advantageous to
follow strictly the analogy of the terminology commonly adopted (though not approved by
all) in the theory of plane curves, according to which the measure of curvature should be
called simply curvature, but the total curvature, the amplitude. But why not be free in
the choice of words, provided they are not meaningless and not liable to a misleading
interpretation?

The position of a figure on the sphere can be either similar to the position of the
corresponding figure on the curved surface, or opposite (inverse). The former is the case
when two lines going out on the curved surface from the same point in different, but not
opposite directions, are represented on the sphere by lines similarly placed, that is, when
the map of the line to the right is also to the right; the latter is the case when the contrary
holds. We shall distinguish these two cases by the positive or negative \emph{sign} of the
measure of curvature. But evidently this distinction can hold only when on each surface
we choose a definite face on which we suppose the figure to lie. On the auxiliary sphere
we shall use always the exterior face, that is, that turned away from the centre; on the
curved surface also there may be taken for the exterior face the one already considered,
or rather that face from which the normal is supposed to be drawn. For, evidently, there
is no change in regard to the similitude of the figures, if on the curved surface both the
figure and the normal be transferred to the opposite side, so long as the image itself
is represented on the same side of the sphere.

The positive or negative sign, which we assign to the \emph{measure} of curvature according
to the position of the infinitely small figure, we extend also to the integral curvature
of a finite figure on the curved surface. However, if we wish to discuss the general case,
some explanations will be necessary, which we can only touch here briefly. So long
as the figure on the curved surface is such that to \emph{distinct} points on itself there correspond
distinct points on the sphere, the definition needs no further explanation. But
whenever this condition is not satisfied, it will be necessary to take into account twice
or several times certain parts of the figure on the sphere. Whence for a similar, or
\PageSep{11}
inverse position, may arise an accumulation of areas, or the areas may partially or
wholly destroy each other. In such a case, the simplest way is to suppose the curved
surface divided into parts, such that each part, considered separately, satisfies the above
condition; to assign to each of the parts its integral curvature, determining this magnitude
by the area of the corresponding figure on the sphere, and the sign by the position
of this figure; and, finally, to assign to the total figure the integral curvature
arising from the addition of the integral curvatures which correspond to the single parts.
So, generally, the integral curvature of a figure is equal to $\int k\, d\sigma$, $d\sigma$~denoting the
element of area of the figure, and $k$~the measure of curvature at any point. The principal
points concerning the geometric representation of this integral reduce to the following.
To the perimeter of the figure on the curved surface (under the restriction
of \Art{3}) will correspond always a closed line on the sphere. If the latter nowhere
intersect itself, it will divide the whole surface of the sphere into two parts, one of
which will correspond to the figure on the curved surface; and its area (taken as
positive or negative according as, with respect to its perimeter, its position is similar,
or inverse, to the position of the figure on the curved surface) will represent the integral
curvature of the figure on the curved surface. But whenever this line intersects
itself once or several times, it will give a complicated figure, to which, however, it is
possible to assign a definite area as legitimately as in the case of a figure without
nodes; and this area, properly interpreted, will give always an exact value for the
integral curvature. However, we must reserve for another occasion\Note{9} the more extended
exposition of the theory of these figures viewed from this very general standpoint.


\Article{7.}

We shall now find a formula which will express the measure of curvature for
any point of a curved surface. Let $d\sigma$~denote the area of an element of this surface;
then $Z\, d\sigma$~will be the area of the projection of this element on the plane of the coordinates
$x$,~$y$; and consequently, if $d\Sigma$~is the area of the corresponding element on the
sphere, $Z\, d\Sigma$~will be the area of its projection on the same plane. The positive or
negative sign of~$Z$ will, in fact, indicate that the position of the projection is similar or
inverse to that of the projected element. Evidently these projections have the same
ratio as to quantity and the same relation as to position as the elements themselves.
Let us consider now a triangular element on the curved surface, and let us suppose
that the coordinates of the three points which form its projection are
\begin{alignat*}{3}
&x,           && y\Add{,} \\
&x + dx,\quad && y + dy\Add{,} \\
&x + \delta x,\quad && y + \delta y\Add{.}
\end{alignat*}
\PageSep{12}
The double area of this triangle will be expressed by the formula
\[
dx·\delta y - dy·\delta x\Add{,}
\]
and this will be in a positive or negative form according as the position of the side
from the first point to the third, with respect to the side from the first point to the
second, is similar or opposite to the position of the $y$-axis of coordinates with respect
to the $x$-axis of coordinates.

In like manner, if the coordinates of the three points which form the projection of
the corresponding element on the sphere, from the centre of the sphere as origin, are
\begin{alignat*}{3}
&X,           && Y\Add{,} \\
&X + dX,\quad && Y + dY\Add{,} \\
&X + \delta X,\quad && Y + \delta Y\Add{,}
\end{alignat*}
the double area of this projection will be expressed by
\[
dX·\delta Y - dY·\delta X\Add{,}
\]
and the sign of this expression is determined in the same manner as above. Wherefore
the measure of curvature at this point of the curved surface will be
\[
k = \frac{dX·\delta Y - dY·\delta X}{dx·\delta y - dy·\delta x}\Add{.}
\]
If now we suppose the nature of the curved surface to be defined according to the third
method considered in \Art{4}, $X$~and~$Y$ will be in the form of functions of the quantities
$x$,~$y$. We shall have, therefore,\Note{10}
\begin{alignat*}{2}
dX &= \frac{\dd X}{\dd x}\, dx &&+ \frac{\dd X}{\dd y}\, dy\Add{,} \\
\delta X &= \frac{\dd X}{\dd x}\, \delta x
          &&+ \frac{\dd X}{\dd y}\, \delta y\Add{,} \\
dY &= \frac{\dd Y}{\dd x}\, dx &&+ \frac{\dd Y}{\dd y}\, dy\Add{,} \\
\delta Y &= \frac{\dd Y}{\dd x}\, \delta x
          &&+ \frac{\dd Y}{\dd y}\, \delta y\Add{.}
\end{alignat*}
When these values have been substituted, the above expression becomes
\[
k = \frac{\dd X}{\dd x}·\frac{\dd Y}{\dd y}
  - \frac{\dd X}{\dd y}·\frac{\dd Y}{\dd x}\Add{.}
\]
\PageSep{13}
Setting, as above,
\[
\frac{\dd z}{\dd x} = t,\quad \frac{\dd z}{\dd y} = u
\]
and also
\[
\frac{\dd^{2} z}{\dd x^{2}} = T,\quad
\frac{\dd^{2} z}{\dd x·\dd y} = U,\quad
\frac{\dd^{2} z}{\dd y^{2}} = V\Add{,}
\]
or
\[
dt = T\, dx + U\, dy,\quad
du = U\, dx + V\, dy\Add{,}
\]
we have from the formulæ given above
\[
X = -tZ,\quad Y = -uZ,\quad (1 - t^{2} - u^{2})Z^{2} = 1\Add{;}
\]
and hence
\begin{gather*}
\begin{alignedat}{2}
dX &= -Z\, dt &&- t\, dZ\Add{,} \\
dY &= -Z\, du &&- u\, dZ\Add{,}
\end{alignedat} \\
(1 + t^{2} + u^{2})\, dZ + Z(t\, dt + u\, du) = 0\Add{;}
\end{gather*}
or\Note{11}
\begin{align*}
dZ &= -Z^{3}(t\, dt + u\, du)\Add{,} \\
dX &= -Z^{3}(1 + u^{2})\, dt + Z^{3} tu\, du\Add{,} \\
dY &= +Z^{3}tu\, dt - Z^{3}(1 + t^{2})\, du\Add{;}\NoteMark
\end{align*}
and so
\begin{align*}
\frac{\dd X}{\dd x} &= Z^{3}\bigl(-(1 + u^{2})T + tuU\bigr)\Add{,} \\
\frac{\dd X}{\dd y} &= Z^{3}\bigl(-(1 + u^{2})U + tuV\bigr)\Add{,} \\
\frac{\dd Y}{\dd x} &= Z^{3}\bigl( tuT - (1 + t^{2})U\bigr)\Add{,} \\
\frac{\dd Y}{\dd y} &= Z^{3}\bigl( tuU - (1 + t^{2})V\bigr)\Add{.}
\end{align*}
Substituting these values in the above expression, it becomes
\begin{align*}
k &= Z^{6}(TV - U^{2}) (1 + t^{2} + u^{2}) = Z^{4} (TV - U^{2}) \\
  &= \frac{TV - U^{2}}{(1 + t^{2} + u^{2})^{2}}\Add{.}
\end{align*}


\Article{8.}

By a suitable choice of origin and axes of coordinates, we can easily make the
values of the quantities $t$,~$u$,~$U$ vanish for a definite point~$A$. Indeed, the first two
\PageSep{14}
conditions will be fulfilled at once if the tangent plane at this point be taken for the
$xy$-plane. If, further, the origin is placed at the point $A$~itself, the expression for
the coordinate~$z$ evidently takes the form
\[
z = \tfrac{1}{2}T°x^{2} + U°xy + \tfrac{1}{2}V°y^{2} + \Omega\Add{,}
\]
where $\Omega$~will be of higher degree than the second. Turning now the axes of $x$~and~$y$
through an angle~$M$ such that
\[
\tan 2M = \frac{2U°}{T° - V°}\Add{,}
\]
it is easily seen that there must result an equation of the form
\[
z = \tfrac{1}{2}Tx^{2} + \tfrac{1}{2}Vy^{2} + \Omega\Add{.}
\]
In this way the third condition is also satisfied. When this has been done, it is evident
that

\Par{I.} If the curved surface be cut by a plane passing through the normal itself and
through the $x$-axis, a plane curve will be obtained, the radius of curvature of which
at the point~$A$ will be equal to~$\dfrac{1}{T}$, the positive or negative sign indicating that the
curve is concave or convex toward that region toward which the coordinates~$z$ are
positive.

\Par{II.} In like manner $\dfrac{1}{V}$~will be the radius of curvature at the point~$A$ of the plane
curve which is the intersection of the surface and the plane through the $y$-axis and
the $z$-axis.

\Par{III.} Setting $z = r \cos\phi$, $y = r \sin \phi$, the equation becomes
\[
z = \tfrac{1}{2}(T\cos^{2}\phi + V\sin^{2}\phi) r^{2} + \Omega\Add{,}
\]
from which we see that if the section is made by a plane through the normal at~$A$
and making an angle~$\phi$ with the $x$-axis, we shall have a plane curve whose radius of
curvature at the point~$A$ will be
\[
\frac{1}{T\cos^{2}\phi + V\sin^{2}\phi}\Add{.}
\]

\Par{IV.} Therefore, whenever we have $T = V$, the radii of curvature in \emph{all} the normal
planes will be equal. But if $T$~and~$V$ are not equal, it is evident that, since for any
value whatever of the angle~$\phi$, $T\cos^{2}\phi + V\sin^{2}\phi$ falls between $T$~and~$V$, the radii of
curvature in the principal sections considered in I.~and~II. refer to the extreme curvatures;
that is to say, the one to the maximum curvature, the other to the minimum,
\PageSep{15}
if $T$~and~$V$ have the same sign. On the other hand, one has the greatest convex
curvature, the other the greatest concave curvature, if $T$~and~$V$ have opposite signs.
These conclusions contain almost all that the illustrious Euler\Note{12} was the first to prove
on the curvature of curved surfaces.

\Par{V.} The measure of curvature at the point~$A$ on the curved surface takes the
very simple form
\[
k = TV,
\]
whence we have the

\begin{Theorem}
The measure of curvature at any point whatever of the surface is equal to a
fraction whose numerator is unity, and whose denominator is the product of the two extreme
radii of curvature of the sections by normal planes.\
\end{Theorem}

At the same time it is clear that the measure of curvature is positive for concavo-concave
or convexo-convex surfaces (which distinction is not essential), but negative
for concavo-convex surfaces. If the surface consists of parts of each kind, then
on the lines separating the two kinds the measure of curvature ought to vanish. Later
we shall make a detailed study of the nature of curved surfaces for which the measure
of curvature everywhere vanishes.


\Article{9.}

The general formula for the measure of curvature given at the end of \Art{7} is
the most simple of all, since it involves only five elements. We shall arrive at a
more complicated formula, indeed, one involving nine elements, if we wish to use the
first method of representing a curved surface. Keeping the notation of \Art{4}, let us
set also
\begin{align*}
\frac{\dd^{2} W}{\dd x^{2}} &= P', &
\frac{\dd^{2} W}{\dd y^{2}} &= Q', &
\frac{\dd^{2} W}{\dd z^{2}} &= R'\Add{,} \\
\frac{\dd^{2} W}{\dd y·\dd z} &= P'', &
\frac{\dd^{2} W}{\dd x·\dd z} &= Q'', &
\frac{\dd^{2} W}{\dd x·\dd y} &= R''\Add{,}
\end{align*}
so that
\begin{alignat*}{4}
dP &= P'\,  &&dx + R''\, &&dy + Q''\, &&dz\Add{,} \\
dQ &= R''\, &&dx + Q' \, &&dy + P''\, &&dz\Add{,} \\
dR &= Q''\, &&dx + P''\, &&dy + R' \, &&dz\Add{.}
\end{alignat*}
Now since $t = -\dfrac{P}{R}$, we find through differentiation
\[
R^{2}\, dt = -R\, dP + P\, dR
  = (PQ'' - RP')\, dx + (PP'' - RR'')\, dy + (PR' - RQ'')\, dz\Add{,}
\]
\PageSep{16}
or, eliminating~$dz$ by means of the equation
\begin{gather*}
P\, dx + Q\, dy + R\, dz = 0, \\
R^{3}\, dt
  = (-R^{2}P' + 2PRQ'' - P^{2}R')\, dx + (PRP'' + QRQ'' - PQR' - R^{2}R'')\, dy.
\end{gather*}
In like manner we obtain
\[
R^{3}\, du
  = (PRP'' + QRQ'' - PQR' - R^{2}R'')\, dx + (-R^{2}Q' + 2QRP'' - Q^{2}R')\, dy\Add{.}
\]
From this we conclude that
\begin{align*}
R^{3}T &= -R^{2}P' + 2PRQ'' - P^{2}R'\Add{,} \\
R^{3}U &= PRP'' + QRQ'' - PQR' - R^{2}R''\Add{,} \\
R^{3}V &= -R^{2}Q' + 2QRP'' - Q^{2}R'\Add{.}
\end{align*}
Substituting these values in the formula of \Art{7}, we obtain for the measure of curvature~$k$
the following symmetric expression:
\begin{multline*}
(P^{2} + Q^{2} + R^{2})^{2}k
  = P^{2}(Q'R' - P''^{2})
  + Q^{2}(P'R' - Q''^{2})
  + R^{2}(P'Q' - R''^{2}) \\
  + 2QR(Q''R'' - P'P'')
  + 2PR(P''R'' - Q'Q'')
  + 2PQ(P''Q'' - R'R'')\Add{.}
\end{multline*}


\Article{10.}

We obtain a still more complicated formula, indeed, one involving fifteen elements,
if we follow the second general method of defining the nature of a curved surface. It
is, however, very important that we develop this formula also. Retaining the notations
of \Art{4}, let us put also
\begin{align*}
\frac{\dd^{2}x}{\dd p^{2}} &= \alpha, &
\frac{\dd^{2}x}{\dd p·\dd q} &= \alpha', &
\frac{\dd^{2}x}{\dd q^{2}} &= \alpha''\Add{,} \\
%
\frac{\dd^{2}y}{\dd p^{2}} &= \beta, &
\frac{\dd^{2}y}{\dd p·\dd q} &= \beta', &
\frac{\dd^{2}y}{\dd q^{2}} &= \beta''\Add{,} \\
%
\frac{\dd^{2}z}{\dd p^{2}} &= \gamma, &
\frac{\dd^{2}z}{\dd p·\dd q} &= \gamma', &
\frac{\dd^{2}z}{\dd q^{2}} &= \gamma''\Add{;} \\
\end{align*}
and let us put, for brevity,
\begin{align*}
bc' - cb' &= A\Add{,} \\
ca' - ac' &= B\Add{,} \\
ab' - ba' &= C\Add{.}
\end{align*}
First we see that
\[
A\, dx + B\, dy + C\, dz = 0,
\]
or
\[
dz = -\frac{A}{C}\, dx - \frac{B}{C}\, dy.
\]
\PageSep{17}
Thus, inasmuch as $z$~may be regarded as a function of $x$,~$y$, we have
\begin{align*}
\frac{\dd z}{\dd x} &= t = -\frac{A}{C}\Add{,} \\
\frac{\dd z}{\dd y} &= u = -\frac{B}{C}\Add{.}
\end{align*}
Then from the formulæ
\[
dx = a\, dp + a'\, dq,\quad
dy = b\, dp + b'\, dq,
\]
we have
\begin{alignat*}{4}
&C\, dp = &&b'\, &&dx - a'\, &&dy\Add{,} \\
&C\, dq =-&&b\,  &&dx + a\,  &&dy\Add{.}
\end{alignat*}
Thence we obtain for the total differentials of $t$,~$u$
\begin{alignat*}{2}
C^{3}\, dt
  &= \left(A\, \frac{\dd C}{\dd p} - C\, \frac{\dd A}{\dd p}\right)(b'\, dx - a'\, dy)
   + \left(C\, \frac{\dd A}{\dd q} - A\, \frac{\dd C}{\dd q}\right)(b \, dx - a \, dy)\Add{,} \\
%
C^{3}\, du
  &= \left(B\, \frac{\dd C}{\dd p} - C\, \frac{\dd B}{\dd p}\right)(b'\, dx - a'\, dy)
   + \left(C\, \frac{\dd B}{\dd q} - B\, \frac{\dd C}{\dd q}\right)(b \, dx - a \, dy)\Add{.}
\end{alignat*}
If now we substitute in these formulæ
\begin{alignat*}{4}
\frac{\dd A}{\dd p} &= c'\beta  &&+ b\gamma'  &&- c\beta'  &&- b'\gamma\Add{,} \\
\frac{\dd A}{\dd q} &= c'\beta' &&+ b\gamma'' &&- c\beta'' &&- b'\gamma'\Add{,} \\
%
\frac{\dd B}{\dd p} &= a'\gamma  &&+ c\alpha'  &&- a\gamma'  &&- c'\alpha\Add{,} \\
\frac{\dd B}{\dd q} &= a'\gamma' &&+ c\alpha'' &&- a\gamma'' &&- c'\alpha'\Add{,} \\
%
\frac{\dd C}{\dd p} &= b'\alpha  &&+ a\beta'  &&- b\alpha'  &&- a'\beta\Add{,} \\
\frac{\dd C}{\dd q} &= b'\alpha' &&+ a\beta'' &&- b\alpha'' &&- a'\beta'\Add{;}
\end{alignat*}
and if we note that the values of the differentials $dt$,~$du$ thus obtained must be equal,
independently of the differentials $dx$,~$dy$, to the quantities $T\, dx + U\, dy$, $U\, dx + V\, dy$
respectively, we shall find, after some sufficiently obvious transformations,
\begin{align*}
C^{3}T &= \alpha Ab'^{2} + \beta Bb'^{2} + \gamma Cb'^{2} \\
&\quad- 2\alpha' Abb' - 2\beta' Bbb' - 2\gamma' Cbb' \\
&\quad+ \alpha'' Ab^{2} + \beta'' Bb^{2} + \gamma'' Cb^{2}\Add{,} \\
\PageSep{18}
C^{3}U &= -\alpha Aa'b' - \beta Ba'b' - \gamma Ca'b' \\
&\quad+ \alpha' A(ab' + ba') + \beta' B(ab' + ba') + \gamma' C(ab' + ba') \\
&\quad- \alpha'' Aab - \beta'' Bab - \gamma'' Cab\Add{,} \\
%
C^{3}V &= \alpha Aa'^{2} + \beta Ba'^{2} + \gamma Ca'^{2} \\
&\quad- 2\alpha' Aaa' - 2\beta' Baa' - 2\gamma' Caa' \\
&\quad+ \alpha'' Aa^{2} + \beta'' Ba^{2} + \gamma'' Ca^{2}\Add{.}
\end{align*}
Hence, if we put, for the sake of brevity,\Note{13}
\begin{alignat*}{4}
&A\alpha   &&+ B\beta   &&+ C\gamma   &&= D\Add{,}\NoteMark
\Tag{(1)} \\
&A\alpha'  &&+ B\beta'  &&+ C\gamma'  &&= D'\Add{,}
\Tag{(2)} \\
&A\alpha'' &&+ B\beta'' &&+ C\gamma'' &&= D''\Add{,}
\Tag{(3)}
\end{alignat*}
we shall have
\begin{align*}
C^{3}T &=  Db'^{2} - 2D'bb' + D'' b^{2}\Add{,} \\
C^{3}U &= -Da'b' + D'(ab' + ba') - D''ab\Add{,} \\
C^{3}V &=  Da'^{2} - 2D'aa' + D''a^{2}\Add{.}
\end{align*}
From this we find, after the reckoning has been carried out,
\[
C^{6}(TV - U^{2}) = (DD'' - D'^{2}) (ab' - ba')^{2} = (DD'' - D'^{2}) C^{2}\Add{,}
\]
and therefore the formula for the measure of curvature
\[
k = \frac{DD'' - D'^{2}}{(A^{2} + B^{2} + C^{2})^{2}}\Add{.}
\]


\Article{11.}

By means of the formula just found we are going to establish another, which may
be counted among the most productive theorems in the theory of curved surfaces.
Let us introduce the following notation:
\begin{alignat*}{4}
&a^{2} &&+ b^{2} &&+ c^{2} &&= E\Add{,} \\
&aa' &&+ bb' &&+ cc' &&= F\Add{,} \\
&a'^{2} &&+ b'^{2} &&+ c'^{2} &&= G\Add{;}
\end{alignat*}
\begin{alignat*}{7}
&a &&\alpha  &&+b &&\beta  &&+c &&\gamma  &&= m  \Add{,}
\Tag{(4)} \\
&a &&\alpha' &&+b &&\beta' &&+c &&\gamma' &&= m' \Add{,}
\Tag{(5)} \\
&a &&\alpha''&&+b &&\beta''&&+c &&\gamma''&&= m''\Add{;}
\Tag{(6)} \displaybreak[1] \\
%
&a'\,&&\alpha  &&+b'\,&&\beta  &&+c'\,&&\gamma  &&= n  \Add{,}
\Tag{(7)} \\
&a'&&\alpha' &&+b'&&\beta' &&+c'&&\gamma' &&= n' \Add{,}
\Tag{(8)} \\
&a'&&\alpha''&&+b'&&\beta''&&+c'&&\gamma''&&= n''\Add{;}
\Tag{(9)}
\end{alignat*}
\[
A^{2} + B^{2} + C^{2} = EG - F^{2} = \Delta\Add{.}
\]
\PageSep{19}

%[** TN: Added parentheses around equation numbers]
Let us eliminate from the equations (1),~(4),~(7) the quantities $\beta$,~$\gamma$, which is done by
multiplying them by $bc' - cb'$, $b'C - c'B$, $cB - bC$ respectively and adding. In this
way we obtain
\begin{multline*}
\bigl(A(bc' - cb') + a(b'C - c'B) + a'(cB - bC)\bigr)\alpha \\
  = D(bc' - cb') + m(b'C - c'B) + n(cB - bC)\Add{,}
\end{multline*}
an equation which is easily transformed into
\[
AD = \alpha\Delta + a(nF - mG) + a'(mF - nE)\Add{.}
\]
Likewise the elimination of $\alpha$,~$\gamma$ or $\alpha$,~$\beta$ from the same equations gives
\begin{alignat*}{4}
&BD &&= \beta \Delta &&+ b(nF - mG) &&+ b'(mF - nE)\Add{,} \\
&CD &&= \gamma\Delta &&+ c(nF - mG) &&+ c'(mF - nE)\Add{.}
\end{alignat*}
Multiplying these three equations by $\alpha''$,~$\beta''$,~$\gamma''$ respectively and adding, we obtain
\[
DD'' = (\alpha\alpha''+ \beta\beta'' + \gamma\gamma'')\Delta
  + m''(nF - mG) + n''(mF - nE)\Add{.}
\Tag{(10)}
\]

%[** TN: Added parentheses around equation numbers]
If we treat the equations (2),~(5),~(8) in the same way, we obtain
\begin{alignat*}{4}
&AD' &&= \alpha'\Delta &&+ a (n'F - m'G) &&+ a'(m'F - n'E)\Add{,} \\
&BD' &&= \beta' \Delta &&+ b (n'F - m'G) &&+ b'(m'F - n'E)\Add{,} \\
&CD' &&= \gamma'\Delta &&+ c (n'F - m'G) &&+ c'(m'F - n'E)\Add{;}
\end{alignat*}
and after these equations are multiplied by $\alpha'$,~$\beta'$,~$\gamma'$ respectively, addition gives
\[
D'^{2} = (\alpha'^{2} + \beta'^{2} + \gamma'^{2})\Delta
  + m'(n'F - m'G) + n'(m'F - n'E)\Add{.}
\]

A combination of this equation with equation~(10) gives
\begin{multline*}
DD'' - D'^{2} = (\alpha\alpha'' + \beta\beta'' + \gamma\gamma''
  - \alpha'^{2} - \beta'^{2} - \gamma'^{2})\Delta \\
  + E(n'^{2} - nn'') + F(nm'' - 2m'n' + mn'') + G(m'^{2} - mm'')\Add{.}
\end{multline*}
It is clear that we have
\[
\frac{\dd E}{\dd p} = 2m,\
\frac{\dd E}{\dd q} = 2m',\quad
\frac{\dd F}{\dd p} = m' + n,\
\frac{\dd F}{\dd q} = m'' + n',\quad
\frac{\dd G}{\dd p} = 2n',\
\frac{\dd G}{\dd q} = 2n'',
\]
or\Note{14}
\begin{align*}
m   &= \tfrac{1}{2}\, \frac{\dd E}{\dd p}, &
m'  &= \tfrac{1}{2}\, \frac{\dd E}{\dd q}, &
m'' &= \frac{\dd F}{\dd q} - \tfrac{1}{2}\, \frac{\dd G}{\dd p}\Add{,}\NoteMark \\
%
n   &= \frac{\dd F}{\dd p} - \tfrac{1}{2}\, \frac{\dd E}{\dd q}, &
n'  &= \tfrac{1}{2}\, \frac{\dd G}{\dd p}, &
n'' &= \tfrac{1}{2}\, \frac{\dd G}{\dd q}\Add{.}
\end{align*}
Moreover, it is easily shown that we shall have
\begin{align*}
%[** TN: Aligning on equals sign]
\alpha\alpha'' + \beta\beta'' + \gamma\gamma''
  - \alpha'^{2} - \beta'^{2} - \gamma'^{2}
  &= \frac{\dd n}{\dd q} - \frac{\dd n'}{\dd p}
   = \frac{\dd m''}{\dd p} - \frac{\dd m'}{\dd q} \\
  &= -\tfrac{1}{2}·\frac{\dd^{2}E}{\dd q^{2}}
  + \frac{\dd^{2}F}{\dd p·\dd q}
  - \tfrac{1}{2}·\frac{\dd^{2}G}{\dd p^{2}}\Add{.}
\end{align*}
\PageSep{20}
If we substitute these different expressions in the formula for the measure of curvature
derived at the end of the preceding article, we obtain the following formula, which
involves only the quantities $E$,~$F$,~$G$ and their differential quotients of the first and
second orders:
\begin{multline*}
4(EG - F^{2})k
  = E\left(\frac{\dd E}{\dd q}·\frac{\dd G}{\dd q}
       - 2 \frac{\dd F}{\dd p}·\frac{\dd G}{\dd q}
  + \biggl(\frac{\dd G}{\dd p}\biggr)^{2}\right) \\
  + F\left(\frac{\dd E}{\dd p}·\frac{\dd G}{\dd q}
         - \frac{\dd E}{\dd q}·\frac{\dd G}{\dd p}
       - 2 \frac{\dd E}{\dd q}·\frac{\dd F}{\dd q}
       + 4 \frac{\dd \Erratum{E}{F}}{\dd p}·\frac{\dd F}{\dd q}
       - 2 \frac{\dd F}{\dd p}·\frac{\dd G}{\dd p}\right) \\
  + G\left(\frac{\dd E}{\dd p}·\frac{\dd G}{\dd p}
       - 2 \frac{\dd E}{\dd p}·\frac{\dd F}{\dd q}
  + \biggl(\frac{\dd E}{\dd q}\biggr)^{2}\right)
  - 2(EG - F^{2})\left(
      \frac{\dd^{2}E}{\dd q^{2}}
   - 2\frac{\dd^{2}F}{\dd p·\dd q}
    + \frac{\dd^{2}G}{\dd p^{2}}
  \right)\Add{.}
\end{multline*}


\Article{12.}

Since we always have
\[
dx^{2} + dy^{2} + dz^{2} = E\, dp^{2} + 2F\, dp·dq + G\, dq^{2},
\]
it is clear that
\[
\Sqrt{E\, dp^{2} + 2F\, dp·dq + G\, dq^{2}}
\]
is the general expression for the linear element on the curved surface. The analysis
developed in the preceding article thus shows us that for finding the measure of curvature
there is no need of finite formulæ, which express the coordinates $x$,~$y$,~$z$ as
functions of the indeterminates $p$,~$q$; but that the general expression for the magnitude
of any linear element is sufficient. Let us proceed to some applications of this very
important theorem.

Suppose that our surface can be developed upon another surface, curved or plane,
so that to each point of the former surface, determined by the coordinates $x$,~$y$,~$z$, will
correspond a definite point of the latter surface, whose coordinates are $x'$,~$y'$,~$z'$. Evidently
$x'$,~$y'$,~$z'$ can also be regarded as functions of the indeterminates $p$,~$q$, and therefore
for the element $\Sqrt{dx'^{2} + dy'^{2} + dz'^{2}}$ we shall have an expression of the form
\[
\Sqrt{E'\, dp^{2} + 2F'\, dp·dq + G'\, dq^{2}}\Add{,}
\]
where $E'$,~$F'$,~$G'$ also denote functions of $p$,~$q$. But from the very notion of the \emph{development}
of one surface upon another it is clear that the elements corresponding to one
another on the two surfaces are necessarily equal. Therefore we shall have identically
\[
E = E',\quad F = F',\quad G = G'.
\]
Thus the formula of the preceding article leads of itself to the remarkable

\begin{Theorem}
If a curved surface is developed upon any other surface whatever, the
measure of curvature in each point remains unchanged.
\end{Theorem}
\PageSep{21}

Also it is evident that
\begin{Theorem}[]
any finite part whatever of the curved surface will retain the
same integral curvature after development upon another surface.
\end{Theorem}

Surfaces developable upon a plane constitute the particular case to which geometers
have heretofore restricted their attention. Our theory shows at once that the
measure of curvature at every point of such surfaces is equal to zero. Consequently,
if the nature of these surfaces is defined according to the third method, we shall have
at every point
\[
\frac{\dd^{2} z}{\dd x^{2}}·\frac{\dd^{2} z}{\dd y^{2}}
  - \left(\frac{\dd^{2}z}{\dd x·\dd y}\right)^{2} = 0\Add{,}
\]
a criterion which, though indeed known a short time ago, has not, at least to our
knowledge, commonly been demonstrated with as much rigor as is desirable.


\Article{13.}

What we have explained in the preceding article is connected with a particular
method of studying surfaces, a very worthy method which may be thoroughly developed
by geometers. When a surface is regarded, not as the boundary of a solid, but
as a flexible, though not extensible solid, one dimension of which is supposed to
vanish, then the properties of the surface depend in part upon the form to which we
can suppose it reduced, and in part are absolute and remain invariable, whatever may
be the form into which the surface is bent. To these latter properties, the study of
which opens to geometry a new and fertile field, belong the measure of curvature and
the integral curvature, in the sense which we have given to these expressions. To
these belong also the theory of shortest lines, and a great part of what we reserve to
be treated later. From this point of view, a plane surface and a surface developable
on a plane, \eg,~cylindrical surfaces, conical surfaces,~etc., are to be regarded as essentially
identical; and the generic method of defining in a general manner the nature of
the surfaces thus considered is always based upon the formula
\[
\Sqrt{E\, dp^{2} + 2F\, dp·dq + G\, dq^{2}},
\]
which connects the linear element with the two indeterminates $p$,~$q$. But before following
this study further, we must introduce the principles of the theory of shortest
lines on a given curved surface.


\Article{14.}

The nature of a curved line in space is generally given in such a way that the
coordinates $x$,~$y$,~$z$ corresponding to the different points of it are given in the form of
functions of a single variable, which we shall call~$w$. The length of such a line from
\PageSep{22}
an arbitrary initial point to the point whose coordinates are $x$,~$y$,~$z$, is expressed by
the integral
\[
%[** TN: Round outer parentheses in the original]
\int dw·\SQRT{
    \left(\frac{dx}{dw}\right)^{2}
  + \left(\frac{dy}{dw}\right)^{2}
  + \left(\frac{dz}{dw}\right)^{2}}\Add{.}
\]
If we suppose that the position of the line undergoes an infinitely small variation, so
that the coordinates of the different points receive the variations $\delta x$,~$\delta y$,~$\delta z$, the variation
of the whole length becomes
\[
\int \frac{dx·d\, \delta x + dy·d\, \delta y + dz·d\, \delta z}
          {\Sqrt{dx^{2} + dy^{2} + dz^{2}}}\Add{,}
\]
which expression we can change into the form\Note{15}
\begin{multline*}
\frac{dx·\delta x + dy·\delta y + dz·\delta z}
     {\Sqrt{dx^{2} + dy^{2} + dz^{2}}} \\
-\int \Biggl(
    \delta x·d\frac{dx}{\Sqrt{dx^{2} + dy^{2} + dz^{2}}}
  + \delta y·d\frac{dy}{\Sqrt{dx^{2} + dy^{2} + dz^{2}}}
  + \delta z·d\frac{dz}{\Sqrt{dx^{2} + dy^{2} + dz^{2}}}
  \Biggr)\Add{.}\NoteMark
\end{multline*}
We know that, in case the line is to be the shortest between its end points, all that
stands under the integral sign must vanish. Since the line must lie on the given
surface, whose nature is defined by the equation
\[
P\, dx + Q\, dy + R\, dz = 0,
\]
the variations $\delta x$,~$\delta y$,~$\delta z$ also must satisfy the equation
\[
P\, \delta x + Q\, \delta y + R\, \delta z = 0,
\]
and from this it follows at once, according to well-known rules, that the differentials
\[
d\frac{dx}{\Sqrt{dx^{2} + dy^{2} + dz^{2}}},\quad
d\frac{dy}{\Sqrt{dx^{2} + dy^{2} + dz^{2}}},\quad
d\frac{dz}{\Sqrt{dx^{2} + dy^{2} + dz^{2}}}
\]
must be proportional to the quantities $P$,~$Q$,~$R$ respectively. Let $dr$~be the element
of the curved line; $\lambda$~the point on the sphere representing the direction of this element;
$L$~the point on the sphere representing the direction of the normal to the curved
surface; finally, let $\xi$,~$\eta$,~$\zeta$ be the coordinates of the point~$\lambda$, and $X$,~$Y$,~$Z$ be those of
the point~$L$ with reference to the centre of the sphere. We shall then have
\[
dx = \xi\, dr,\quad
dy = \eta\, dr,\quad
dz = \zeta\, dr\Add{,}
\]
from which we see that the above differentials become $d\xi$,~$d\eta$,~$d\zeta$. And since the
quantities $P$,~$Q$,~$R$ are proportional to $X$,~$Y$,~$Z$, the character of shortest lines is
expressed by the equations
\[
\frac{d\xi}{X} = \frac{d\eta}{Y} = \frac{d\zeta}{Z}\Add{.}
\]
\PageSep{23}
Moreover, it is easily seen that
\[
\Sqrt{d\xi^{2} + d\eta^{2} + d\zeta^{2}}
\]
is equal to the small arc on the sphere which measures the angle between the directions
of the tangents at the beginning and at the end of the element~$dr$, and is thus
equal to~$\dfrac{dr}{\rho}$, if $\rho$~denotes the radius of curvature of the shortest line at this point.
Thus we shall have
\[
\rho\, d\xi = X\, dr,\quad
\rho\, d\eta = Y\, dr,\quad
\rho\, d\zeta = Z\, dr\Add{.}
\]


\Article{15.}

Suppose that an infinite number of shortest lines go out from a given point~$A$
on the curved surface, and suppose that we distinguish these lines from one another
by the angle that the first element of each of them makes with the first element of
one of them which we take for the first. Let $\phi$~be that angle, or, more generally, a
function of that angle, and $r$~the length of such a shortest line from the point~$A$ to
the point whose coordinates are $x$,~$y$,~$z$. Since to definite values of the variables $r$,~$\phi$
there correspond definite points of the surface, the coordinates $x$,~$y$,~$z$ can be regarded
as functions of $r$,~$\phi$. We shall retain for the notation $\lambda$, $L$, $\xi$,~$\eta$,~$\zeta$, $X$,~$Y$,~$Z$ the same
meaning as in the preceding article, this notation referring to any point whatever on
any one of the shortest lines.

All the shortest lines that are of the same length~$r$ will end on another line
whose length, measured from an arbitrary initial point, we shall denote by~$v$. Thus $v$~can
be regarded as a function of the indeterminates $r$,~$\phi$, and if $\lambda'$~denotes the point
on the sphere corresponding to the direction of the element~$dv$, and also $\xi'$,~$\eta'$,~$\zeta'$
denote the coordinates of this point with reference to the centre of the sphere, we
shall have
\[
\frac{\dd x}{\dd\phi} = \xi'·\frac{\dd v}{\dd\phi},\quad
\frac{\dd y}{\dd\phi} = \eta'·\frac{\dd v}{\dd\phi},\quad
\frac{\dd z}{\dd\phi} = \zeta'·\frac{\dd v}{\dd\phi}\Add{.}
\]
From these equations and from the equations
\[
\frac{\dd x}{\dd r} = \xi,\quad
\frac{\dd y}{\dd r} = \eta,\quad
\frac{\dd z}{\dd r} = \zeta
\]
we have
\[
\frac{\dd x}{\dd r}·\frac{\dd x}{\dd\phi} +
\frac{\dd y}{\dd r}·\frac{\dd y}{\dd\phi} +
\frac{\dd z}{\dd r}·\frac{\dd z}{\dd\phi}
  = (\xi\xi' + \eta\eta' + \zeta\zeta')·\frac{\dd v}{\dd\phi}
  = \cos \lambda\lambda'·\frac{\dd v}{\dd\phi}\Add{.}
\]
\PageSep{24}
Let $S$~denote the first member of this equation, which will also be a function of $r$,~$\phi$.
Differentiation of~$S$ with respect to~$r$ gives
\begin{align*}
\frac{\dd S}{\dd r}
  &= \frac{\dd^{2} x}{\dd r^{2}}·\frac{\dd x}{\dd\phi}
   + \frac{\dd^{2} y}{\dd r^{2}}·\frac{\dd y}{\dd\phi}
   + \frac{\dd^{2} z}{\dd r^{2}}·\frac{\dd z}{\dd\phi}
   + \tfrac{1}{2}·\frac{\dd\left(
       \biggl(\dfrac{\dd x}{\dd r}\biggr)^{2}
     + \biggl(\dfrac{\dd y}{\dd r}\biggr)^{2}
     + \biggl(\dfrac{\dd z}{\dd r}\biggr)^{2}
   \right)}{\dd \phi} \\
  &= \frac{\dd\xi}{\dd r}·\frac{\dd x}{\dd\phi}
   + \frac{\dd\eta}{\dd r}·\frac{\dd y}{\dd\phi}
   + \frac{\dd\zeta}{\dd r}·\frac{\dd z}{\dd\phi}
   + \tfrac{1}{2}·\frac{\dd(\xi^{2} + \eta^{2} + \zeta^{2})}{\dd\phi}\Add{.}
\end{align*}
But
\[
\xi^{2} + \eta^{2} + \zeta^{2} = 1,
\]
and therefore its differential is equal to zero; and by the preceding article we have,
if $\rho$~denotes the radius of curvature of the line~$r$,
\[
\frac{\dd\xi}{\dd r} = \frac{X}{\rho},\quad
\frac{\dd\eta}{\dd r} = \frac{Y}{\rho},\quad
\frac{\dd\zeta}{\dd r} = \frac{Z}{\rho}\Add{.}
\]
Thus we have
\[
\frac{\dd S}{\dd r}
  = \frac{1}{\rho}·(X\xi' + Y\eta' + Z\zeta')·\frac{\dd v}{\dd\phi}
  = \frac{1}{\rho}·\cos L\lambda'·\frac{\dd v}{\dd\phi} = 0
\]
since $\lambda'$~evidently lies on the great circle whose pole is~$L$. From this we see that
$S$~is independent of~$r$, and is, therefore, a function of $\phi$~alone. But for $r = 0$ we evidently
have $v = 0$, consequently $\dfrac{\dd v}{\dd\phi} = 0$, and $S = 0$ independently of~$\phi$. Thus, in general,
we have necessarily $S = 0$, and so $\cos\lambda\lambda' = 0$, \ie, $\lambda\lambda' = 90°$. From this follows the

\begin{Theorem}
If on a curved surface an infinite number of shortest lines of equal length
be drawn from the same initial point, the lines joining their extremities will be normal to
each of the lines.
\end{Theorem}

We have thought it worth while to deduce this theorem from the fundamental
property of shortest lines; but the truth of the theorem can be made apparent without
any calculation by means of the following reasoning. Let $AB$,~$AB'$ be two
shortest lines of the same length including at~$A$ an infinitely small angle, and let us
suppose that one of the angles made by the element~$BB'$ with the lines $BA$,~$B'A$
differs from a right angle by a finite quantity. Then, by the law of continuity, one
will be greater and the other less than a right angle. Suppose the angle at~$B$ is
equal to~$90° - \omega$, and take on the line~$AB$ a point~$C$, such that
\[
BC = BB'·\cosec \omega.
\]
Then, since the infinitely small triangle~$BB'C$ may be regarded as plane, we shall have
\[
CB' = BC·\cos \omega,
\]
\PageSep{25}
and consequently
\[
AC + CB' = AC + BC·\cos \omega
  = AB - BC·(1- \cos \omega)
  = AB' - BC·(1 - \cos \omega),
\]
\ie, the path from $A$~to~$B'$ through the point~$C$ is shorter than the shortest line,
\QEA


\Article{16.}

%[** TN: In-line theorem, marked non-semantically]
With the theorem of the preceding article we associate another, which we state
as follows: \textit{If on a curved surface we imagine any line whatever, from the different points
of which are drawn at right angles and toward the same side an infinite number of shortest
lines of the same length, the curve which joins their other extremities will cut each of the
lines at right angles.} For the demonstration of this theorem no change need be made
in the preceding analysis, except that $\phi$~must denote the length of the \emph{given} curve
measured from an arbitrary point; or rather, a function of this length. Thus all of
the reasoning will hold here also, with this modification, that $S = 0$ for $r = 0$ is
now implied in the hypothesis itself. Moreover, this theorem is more general than
the preceding one, for we can regard it as including the first one if we take for the
given line the infinitely small circle described about the centre~$A$. Finally, we may
say that here also geometric considerations may take the place of the analysis, which,
however, we shall not take the time to consider here, since they are sufficiently
obvious.


\Article{17.}

We return to the formula
\[
\Sqrt{E\, dp^{2} + 2F\, dp·dq + G\, dq^{2}},
\]
which expresses generally the magnitude of a linear element on the curved surface,
and investigate, first of all, the geometric meaning of the coefficients $E$,~$F$,~$G$. We
have already said in \Art{5} that two systems of lines may be supposed to lie on the
curved surface, $p$~being variable, $q$~constant along each of the lines of the one system;
and $q$~variable, $p$~constant along each of the lines of the other system. Any point
whatever on the surface can be regarded as the intersection of a line of the first
system with a line of the second; and then the element of the first line adjacent to
this point and corresponding to a variation~$dp$ will be equal to~$\sqrt{E}·dp$, and the
element of the second line corresponding to the variation~$dq$ will be equal to~$\sqrt{G}·dq$.
Finally, denoting by~$\omega$ the angle between these elements, it is easily seen that we
shall have
\[
\cos \omega = \frac{F}{\sqrt{EG}}.
\]
\PageSep{26}
Furthermore, the area of the surface element in the form of a parallelogram between
the two lines of the first system, to which correspond $q$,~$q + dq$, and the two lines of
the second system, to which correspond $p$,~$p + dp$, will be
\[
\Sqrt{EG - F^{2}}\, dp·dq.
\]

Any line whatever on the curved surface belonging to neither of the two systems
is determined when $p$~and~$q$ are supposed to be functions of a new variable, or
one of them is supposed to be a function of the other. Let $s$~be the length of such
a curve, measured from an arbitrary initial point, and in either direction chosen as
positive. Let $\theta$~denote the angle which the element
\[
ds = \Sqrt{E\, dp^{2} + 2F\, dp·dq + G\, dq^{2}}
\]
makes with the line of the first system drawn through the initial point of the element,
and, in order that no ambiguity may arise, let us suppose that this angle is
measured from that branch of the first line on which the values of~$p$ increase, and is
taken as positive toward that side toward which the values of~$q$ increase. These conventions
being made, it is easily seen that
\begin{align*}
\cos \theta·ds &= \sqrt{E}·dp + \sqrt{G}·\cos \omega·dq
  = \frac{E\, dp + F\, dq}{\sqrt{E}}\Add{,} \\
\sin \theta·ds &= \sqrt{G}·\sin \omega·dq
  = \frac{\sqrt{(EG - F^{2})}·dq}{\sqrt{E}}\Add{.}
\end{align*}


\Article{18.}

We shall now investigate the condition that this line be a shortest line. Since
its length~$s$ is expressed by the integral
\[
s = \int \Sqrt{E\, dp^{2} + 2F\, dp·dq + G\, dq^{2}}\Add{,}
\]
the condition for a minimum requires that the variation of this integral arising from
an infinitely small change in the position become equal to zero. The calculation, for
our purpose, is more simply made in this case, if we regard $p$ as a function of~$q$.
When this is done, if the variation is denoted by the characteristic~$\delta$, we have
\begin{align*}
\delta s &= \int \frac{\left(
    \dfrac{\dd E}{\dd p}·dp^{2}
  + 2\dfrac{\dd F}{\dd p}·dp·dq
  + \dfrac{\dd G}{\dd p}·dq^{2}
  \right) \delta p
  + (2E\, dp + 2F\, dq)\, d\, \delta p}{2\, ds} \displaybreak[1] \\
  &= \frac{E\, dp + F\, dq}{ds}·\delta p \\
\PageSep{27}
  &\qquad+ \int \delta p \left(\frac{
    \dfrac{\dd E}{\dd p}·dp^{2}
  + 2\dfrac{\dd F}{\dd p}·dp·dq
  + \dfrac{\dd G}{\dd p}·dq^{2}}{2\, ds}
  - d·\frac{E\, dp + F\, dq}{ds}\right)
\end{align*}
and we know that what is included under the integral sign must vanish independently
of~$\delta p$. Thus we have
\begin{multline*}
\frac{\dd E}{\dd p}·dp^{2}
  + 2\frac{\dd F}{\dd p}·dp·dq
  + \frac{\dd G}{\dd p}·dq^{2}
  = 2\, ds·d·\frac{E\, dp + F\, dq}{ds} \\
\begin{aligned}
&= 2\, ds·d·\left(\sqrt{E}·\cos\theta\right) \\ %[** TN: Added parentheses]
&= \frac{ds·dE·\cos\theta}{\sqrt{E}}
 - 2\, ds·d\theta·\sqrt{E}·\sin\theta \\
%** Translator's note from corrigenda: The original and the Latin reprints **
%** lack the factor 2; the correction is made in all the translations.     **
&= \frac{(E\, dp + F\, dq)\, dE}{E} - \Erratum{}{2}\Sqrt{EG - F^{2}}·\Erratum{dp}{dq}·d\theta \\
&= \left(\frac{E\, dp + F\, dq}{E}\right)
 ·\left(\frac{\dd E}{\dd p}·dp + \frac{\dd E}{\dd q}·dq\right)
  - 2\Sqrt{EG - F^{2}}·dq·d\theta\Add{.}
\end{aligned}
\end{multline*}
This gives the following conditional equation for a shortest line:
\begin{multline*}
\Sqrt{EG - F^{2}}·d\theta
  = \frac{1}{2}·\frac{F}{E}·\frac{\dd E}{\dd p}·dp
  + \frac{1}{2}·\frac{F}{E}·\frac{\dd E}{\dd q}·dq
  + \frac{1}{2}·\frac{\dd E}{\dd q}·dp \\
  - \frac{\dd F}{\dd p}·dp - \frac{1}{2}·\frac{\dd G}{\dd p}·dq\Add{,}
\end{multline*}
which can also be written
\[
\Sqrt{EG - F^{2}}·d\theta
  = \frac{1}{2}·\frac{F}{E}·dE
  + \frac{1}{2}·\frac{\dd E}{\dd q}·dp
  - \frac{\dd F}{\dd p}·dp
  - \frac{1}{2}·\frac{\dd G}{\dd p}·dq\Add{.}
\]
From this equation, by means of the equation
\[
\cot\theta = \frac{E}{\Sqrt{EG - F^{2}}}·\frac{dp}{dq}
  + \frac{F}{\Sqrt{EG - F^{2}}}\Add{,}
\]
it is also possible to eliminate the angle~$\theta$, and to derive a differential equation of
the second order between $p$~and~$q$, which, however, would become more complicated
and less useful for applications than the preceding.


\Article{19.}

The general formulæ, which we have derived in \Arts{11}{18} for the measure of
curvature and the variation in the direction of a shortest line, become much simpler
if the quantities $p$,~$q$ are so chosen that the lines of the first system cut everywhere
\PageSep{28}
orthogonally the lines of the second system; \ie, in such a way that we have generally
$\omega = 90°$, or $F = 0$. Then the formula for the measure of curvature becomes
\[
4E^{2}G^{2}k
  = E·\frac{\dd E}{\dd q}·\frac{\dd G}{\dd q}
  + E\left(\frac{\dd G}{\dd p}\right)^{2}
  + G·\frac{\dd E}{\dd p}·\frac{\dd G}{\dd p}
  + G\left(\frac{\dd E}{\dd q}\right)^{2}
  - 2EG\left(\frac{\dd^{2} E}{\dd q^{2}} + \frac{\dd^{2} G}{\dd p^{2}}\right),
\]
and for the variation of the angle~$\theta$
\[
\sqrt{EG}·d\theta
  = \frac{1}{2}·\frac{\dd E}{\dd q}·dp
  - \frac{1}{2}·\frac{\dd G}{\dd p}·dq\Add{.}
\]

Among the various cases in which we have this condition of orthogonality, the
most important is that in which all the lines of one of the two systems, \eg, the
first, are shortest lines. Here for a constant value of~$q$ the angle~$\theta$ becomes equal to
zero, and therefore the equation for the variation of~$\theta$ just given shows that we must
have $\dfrac{\dd E}{\dd q} = 0$, or that the coefficient~$E$ must be independent of~$\Erratum{g}{q}$; \ie, $E$~must be
either a constant or a function of $p$~alone. It will be simplest to take for~$p$
the length of each line of the first system, which length, when all the lines of the
first system meet in a point, is to be measured from this point, or, if there is no
common intersection, from any line whatever of the second system. Having made
these conventions, it is evident that $p$~and~$q$ denote now the same quantities that
were expressed in \Arts{15}{16} by $r$~and~$\phi$, and that $E = 1$. Thus the two preceding
formulæ become:
\begin{align*}
4G^{2}k
  &= \left(\frac{\dd G}{\dd p}\right)^{2} - 2G\, \frac{\dd^{2} G}{\dd p^{2}}\Add{,} \\
\sqrt{G}·d\theta
  &= -\frac{1}{2}·\frac{\dd G}{\dd p}·dq\Add{;}
\end{align*}
or, setting $\sqrt{G} = m$,
\[
k = -\frac{1}{m}·\frac{\dd^{2} m}{\dd p^{2}},\quad
d\theta = -\frac{\dd m}{\dd p}·dq\Add{.}
\]
Generally speaking, $m$~will be a function of $p$,~$q$, and $m\, dq$~the expression for the element
of any line whatever of the second system. But in the particular case where
all the lines~$p$ go out from the same point, evidently we must have $m = 0$ for $p = 0$.
Furthermore, in the case under discussion we will take for~$q$ the angle itself which
the first element of any line whatever of the first system makes with the element of
any one of the lines chosen arbitrarily. Then, since for an infinitely small value of~$p$
the element of a line of the second system (which can be regarded as a circle
described with radius~$p$) is equal to~$p\, dq$, we shall have for an infinitely small value
of~$p$, $m = p$, and consequently, for $p = 0$, $m = 0$ at the same time, and $\dfrac{\dd m}{\dd p} = 1$.
\PageSep{29}


\Article{20.}

We pause to investigate the case in which we suppose that $p$~denotes in a general
manner the length of the shortest line drawn from a fixed point~$A$ to any other
point whatever of the surface, and $q$~the angle that the first element of this line
makes with the first element of another given shortest line going out from~$A$. Let
$B$ be a definite point in the latter line, for which $q = 0$, and $C$~another definite point
of the surface, at which we denote the value of~$q$ simply by~$A$. Let us suppose the
points $B$,~$C$ joined by a shortest line, the parts of which, measured from~$B$, we denote
in a general way, as in \Art{18}, by~$s$; and, as in the same article, let us denote by~$\theta$
the angle which any element~$ds$ makes with the element~$dp$; finally, let us denote
by $\theta°$,~$\theta'$ the values of the angle~$\theta$ at the points $B$,~$C$. We have thus on the curved
surface a triangle formed by shortest lines. The angles of this triangle at $B$~and~$C$
we shall denote simply by the same letters, and $B$~will be equal to~$180° - \theta$, $C$~to $\theta'$~itself.
But, since it is easily seen from our analysis that all the angles are supposed
to be expressed, not in degrees, but by numbers, in such a way that the angle $57°\, 17'\, 45''$,
to which corresponds an arc equal to the radius, is taken for the unit, we must set
\[
\theta° = \pi - B,\quad \theta' = C\Add{,}
\]
where $2\pi$~denotes the circumference of the sphere. Let us now examine the integral
curvature of this triangle, which is equal to
\[
\int k\, d\sigma,
\]
$d\sigma$~denoting a surface element of the triangle. Wherefore, since this element is expressed
by~$m\, dp·dq$, we must extend the integral
\[
\iint \Typo{}{k}m\, dp·dq
\]
over the whole surface of the triangle. Let us begin by integration with respect to~$p$,
which, because
\[
k = -\frac{1}{m}·\frac{\dd^{2} m}{\dd p^{2}},
\]
gives
\[
dq·\left(\text{const.} - \frac{\dd m}{\dd p}\right),
\]
for the integral curvature of the area lying between the lines of the first system, to
which correspond the values $q$,~$q + dq$ of the second indeterminate. Since this integral
\PageSep{30}
curvature must vanish for $p = 0$, the constant introduced by integration must be
equal to the value of~$\dfrac{\dd m}{\dd q}$ for $p = 0$, \ie,~equal to unity. Thus we have
\[
dq\left(1 - \frac{\dd m}{\dd p}\right),
\]
where for $\dfrac{\dd m}{\dd p}$ must be taken the value corresponding to the end of this area on the
line~$CB$. But on this line we have, by the preceding article,
\[
\frac{\dd m}{\dd q}·dq = -d\theta,
\]
whence our expression is changed into $dq + d\theta$. Now by a second integration, taken
from $q = 0$ to $q = A$, we obtain for the integral curvature
\[
A + \theta'- \theta°,
\]
or
\[
A + B + C - \pi.
\]

The integral curvature is equal to the area of that part of the sphere which corresponds
to the triangle, taken with the positive or negative sign according as the
curved surface on which the triangle lies is concavo-concave or concavo-convex. For
unit area will be taken the square whose side is equal to unity (the radius of the
sphere), and then the whole surface of the sphere becomes equal to~$4\pi$. Thus the
part of the surface of the sphere corresponding to the triangle is to the whole surface
of the sphere as $±(A + B + C - \pi)$ is to~$4\pi$. This theorem, which, if we mistake
not, ought to be counted among the most elegant in the theory of curved surfaces,
may also be stated as follows:

\begin{Theorem}[]
The excess over~$180°$ of the sum of the angles of a triangle formed by shortest lines
on a concavo-concave curved surface, or the deficit from~$180°$ of the sum of the angles of
a triangle formed hy shortest lines on a concavo-convex curved surface, is measured by the
area of the part of the sphere which corresponds, through the directions of the normals, to
that triangle, if the whole surface of the sphere is set equal to $720$~degrees.
\end{Theorem}

More generally, in any polygon whatever of $n$~sides, each formed by a shortest
line, the excess of the sum of the angles over $(2n - 4)$~right angles, or the deficit from
$(2n - 4)$~right angles (according to the nature of the curved surface), is equal to the
area of the corresponding polygon on the sphere, if the whole surface of the sphere is
set equal to $720$~degrees. This follows at once from the preceding theorem by dividing
the polygon into triangles.
\PageSep{31}


\Article{21.}

Let us again give to the symbols $p$,~$q$, $E$,~$F$,~$G$, $\omega$ the general meanings which
were given to them above, and let us further suppose that the nature of the curved
surface is defined in a similar way by two other variables, $p'$,~$q'$, in which case the
general linear element is expressed by
\[
\Sqrt{E'\, dp'^{2} + 2F'\, dp'·dq' + G'\, dq'^{2}}\Add{.}
\]
Thus to any point whatever lying on the surface and defined by definite values of
the variables $p$,~$q$ will correspond definite values of the variables $p'$,~$q'$, which will
therefore be functions of $p$,~$q$. Let us suppose we obtain by differentiating them
\begin{alignat*}{2}
dp' &= \alpha\, dp &&+ \beta\,  dq\Add{,} \\
dq' &= \gamma\, dp &&+ \delta\, dq\Add{.}
\end{alignat*}
We shall now investigate the geometric meaning of the coefficients $\alpha$,~$\beta$,~$\gamma$,~$\delta$.

Now \emph{four} systems of lines may thus be supposed to lie upon the curved surface,
for which $p$,~$q$, $p'$,~$q'$ respectively are constants. If through the definite point to
which correspond the values $p$,~$q$, $p'$,~$q'$ of the variables we suppose the four lines
belonging to these different systems to be drawn, the elements of these lines, corresponding
to the positive increments $dp$,~$dq$, $dp'$,~$dq'$, will be
\[
\sqrt{E}·dp,\quad
\sqrt{G}·dq,\quad
\sqrt{E'}·dp',\quad
\sqrt{G'}·dq'.
\]
The angles which the directions of these elements make with an arbitrary fixed direction
we shall denote by $M$,~$N$, $M'$,~$N'$, measuring them in the sense in which the
second is placed with respect to the first, so that $\sin(N - M)$ is positive. Let us
suppose (which is permissible) that the fourth is placed in the same sense with respect
to the third, so that $\sin(N' - M')$ also is positive. Having made these conventions,
if we consider another point at an infinitely small distance from the first point, and
to which correspond the values $p + dp$, $q + dq$, $p' + dp'$, $q' + dq'$ of the variables, we
see without much difficulty that we shall have generally, \ie, independently of the
values of the increments $dp$,~$dq$, $dp'$,~$dq'$,
\[
\sqrt{E}·dp·\sin M + \sqrt{G}·dq·\sin N
 = \sqrt{E'}·dp'·\sin M' + \sqrt{G'}·dq'·\sin N'\Add{,}
\]
since each of these expressions is merely the distance of the new point from the line
from which the angles of the directions begin. But we have, by the notation introduced
above,
\[
N - M = \omega.
\]
In like manner we set
\[
N' - M' = \omega',
\]
\PageSep{32}
and also
\[
N - M' = \psi.
\]
Then the equation just found can be thrown into the following form:
\begin{multline*}
\sqrt{E}·dp · \sin(M' - \omega + \psi) + \sqrt{G}·dq·\sin(M' + \psi) \\
= \sqrt{E'}·dp'·\sin M' + \sqrt{G'}·dq'·\sin(M' + \omega')\Add{,}
\end{multline*}
or
\begin{multline*}
\sqrt{E}·dp·\sin(N' - \omega - \omega' + \psi)
  - \sqrt{G}·dq·\sin(N' - \omega' + \psi) \\
  = \sqrt{E'}·dp'·\sin(N' - \omega') + \sqrt{G'}·dq'·\sin N'\Add{.}
\end{multline*}
And since the equation evidently must be independent of the initial direction, this
direction can be chosen arbitrarily. Then, setting in the second formula $N' = 0$, or in
the first $M' = 0$, we obtain the following equations:
\begin{align*}
\sqrt{E'}·\sin \omega'·dp'
  &= \sqrt{E}·\sin(\omega + \omega' - \psi)·dp
   + \sqrt{G}·\sin(\omega' - \psi)·dq\Add{,} \\
\sqrt{G'}·\sin \omega'·dq'
  &= \sqrt{E}·\sin(\psi - \omega)·dp + \sqrt{G}·\sin\psi·dq\Add{;}
\end{align*}
and these equations, since they must be identical with
\begin{alignat*}{2}
dp' &= \alpha\, dp &&+ \beta\,  dq\Add{,} \\
dq' &= \gamma\, dp &&+ \delta\, dq\Add{,}
\end{alignat*}
determine the coefficients $\alpha$,~$\beta$,~$\gamma$,~$\delta$. We shall have
\begin{align*}
\alpha &= \sqrt{\frac{E}{E'}}·\frac{\sin(\omega + \omega' - \psi)}{\sin\omega'}, &
\beta  &= \sqrt{\frac{G}{E'}}·\frac{\sin(\omega' - \psi)}{\sin\omega'}\Add{,} \\
\gamma &= \sqrt{\frac{E}{G'}}·\frac{\sin(\psi - \omega)}{\sin\omega'}, &
\delta &= \sqrt{\frac{G}{G'}}·\frac{\sin\psi}{\sin\omega'}\Add{.}
\end{align*}
These four equations, taken in connection with the equations
\begin{align*}
\cos\omega  &= \frac{F}{\sqrt{EG}}, &
\cos\omega' &= \frac{F'}{\sqrt{E'G'}}, \\
\sin\omega  &= \sqrt{\frac{EG - F^{2}}{EG}}, &
\sin\omega' &= \sqrt{\frac{E'G' - F'^{2}}{E'G'}},
\end{align*}
may be written
\begin{align*}
\alpha\Sqrt{E'G' - F'^{2}} &= \sqrt{EG'}·\sin(\omega + \omega' - \psi)\Add{,} \\
\beta \Sqrt{E'G' - F'^{2}} &= \sqrt{GG'}·\sin(\omega' - \psi)\Add{,} \\
\gamma\Sqrt{E'G' - F'^{2}} &= \sqrt{EE'}·\sin(\psi - \omega)\Add{,} \\
\delta\Sqrt{E'G' - F'^{2}} &= \sqrt{GE'}·\sin \psi\Add{.}
\end{align*}

Since by the substitutions
\begin{alignat*}{2}
dp' &= \alpha\, dp &&+ \beta\,  dq, \\
dq' &= \gamma\, dp &&+ \delta\, dq\Add{,}
\end{alignat*}
\PageSep{33}
the trinomial
\[
E'\, dp'^{2} + 2F'\, dp'·dq' + G'\, dq'^{2}
\]
is transformed into
\[
E\, dp^{2} + 2F\, dp·dq + G\, dq^{2},
\]
we easily obtain
\[
EG - F^{2} = (E'G' - F'^{2})(\alpha\delta - \beta\gamma)^{2}\Add{;}
\]
and since, \textit{vice versa}, the latter trinomial must be transformed into the former by the
substitution
\[
(\alpha\delta - \beta\gamma)\, dp =  \delta\, dp' - \beta\, dq',\quad
(\alpha\delta - \beta\gamma)\, dq = -\gamma\, dp' + \alpha\, dq',
\]
we find\Note{16}
\begin{align*}
E\delta^{2} - 2F\gamma\delta + G\gamma^{2}
  &= \frac{EG - F^{2}}{E'G' - F'^{2}}·E'\Add{,} \\
%
-E\beta\delta + F(\alpha\delta + \beta\gamma) - G\alpha\gamma
  &= \frac{EG - F^{2}}{E'G' - F'^{2}}·F'\Add{,}\NoteMark \\
%
E\beta^{2} - 2F\alpha\beta + G\alpha^{2}
  &= \frac{EG - F^{2}}{E'G' - F'^{2}}·G'\Add{.}
\end{align*}


\Article{22.}

From the general discussion of the preceding article we proceed to the very
extended application in which, while keeping for $p$,~$q$ their most general meaning, we
take for $p'$,~$q'$ the quantities denoted in \Art{15} by $r$,~$\phi$. We shall use $r$,~$\phi$ here
also in such a way that, for any point whatever on the surface, $r$~will be the shortest
distance from a fixed point, and $\phi$~the angle at this point between the first element
of~$r$ and a fixed direction. We have thus
\[
E' = 1,\quad
F' = 0,\quad
\omega' = 90°.
\]
Let us set also
\[
\sqrt{G'} = m,
\]
so that any linear element whatever becomes equal to
\[
\Sqrt{dr^{2} + m^{2}\, d\phi^{2}}.
\]
Consequently, the four equations deduced in the preceding article for $\alpha$,~$\beta$,~$\gamma$,~$\delta$ give
\begin{align*}
\sqrt{E}·\cos(\omega - \psi) = \frac{\dd r}{\dd p}\Add{,}
\Tag{(1)} \\
\sqrt{G}·\cos \psi = \frac{\dd r}{\dd q}\Add{,}
\Tag{(2)} \displaybreak[1] \\
\PageSep{34}
\sqrt{E}·\sin(\psi - \omega) = m·\frac{\dd\phi}{\dd p}\Add{,}
\Tag{(3)} \\
\sqrt{G}·\sin\psi = m·\frac{\dd\phi}{\dd q}\Add{.}
\Tag{(4)}
\end{align*}
But the last and the next to the last equations of the preceding article give
\begin{gather*}
EG - F^{2}
  = E\left(\frac{\dd r}{\dd q}\right)^{2}
  - 2F·\frac{\dd r}{\dd p}·\frac{\dd r}{\dd q}
  + G\left(\frac{\dd r}{\dd p}\right)^{2}\Add{,}
\Tag{(5)} \\
\left(E·\frac{\dd r}{\dd q} - F·\frac{\dd r}{\dd p}\right)·\frac{\dd\phi}{\dd q}
  = \left(F·\frac{\dd r}{\dd q} - G·\frac{\dd r}{\dd p}\right)·\frac{\dd\phi}{\dd p}\Add{.}
\Tag{(6)}
\end{gather*}

From these equations must be determined the quantities $r$,~$\phi$,~$\psi$ and (if need be)~$m$,
as functions of $p$~and~$q$. Indeed, integration of equation~(5) will give~$r$; $r$~being
found, integration of equation~(6) will give~$\phi$; and one or other of equations (1),~(2)
will give $\psi$~itself. Finally, $m$~is obtained from one or other of equations (3),~(4).

The general integration of equations (5),~(6) must necessarily introduce two arbitrary
functions. We shall easily understand what their meaning is, if we remember
that these equations are not limited to the case we are here considering, but are
equally valid if $r$~and~$\phi$ are taken in the more general sense of \Art{16}, so that $r$~is
the length of the shortest line drawn normal to a fixed but arbitrary line, and $\phi$~is
an arbitrary function of the length of that part of the fixed line which is intercepted
between any shortest line and an arbitrary fixed point. The general solution must
embrace all this in a general way, and the arbitrary functions must go over into
definite functions only when the arbitrary line and the arbitrary functions of its
parts, which $\phi$~must represent, are themselves defined. In our case an infinitely
small circle may be taken, having its centre at the point from which the distances~$r$
are measured, and $\phi$~will denote the parts themselves of this circle, divided by the
radius. Whence it is easily seen that the equations (5),~(6) are quite sufficient for
our case, provided that the functions which they leave undefined satisfy the condition
which $r$~and~$\phi$ satisfy for the initial point and for points at an infinitely small
distance from this point.

Moreover, in regard to the integration itself of the equations (5),~(6), we know
that it can be reduced to the integration of ordinary differential equations, which, however,
often happen to be so complicated that there is little to be gained by the reduction.
On the contrary, the development in series, which are abundantly sufficient for
practical requirements, when only a finite portion of the surface is under consideration,
presents no difficulty; and the formulæ thus derived open a fruitful source for
\PageSep{35}
the solution of many important problems. But here we shall develop only a single
example in order to show the nature of the method.


\Article{23.}

We shall now consider the case where all the lines for which $p$~is constant are
shortest lines cutting orthogonally the line for which $\phi = 0$, which line we can regard
as the axis of abscissas. Let $A$~be the point for which $r = 0$, $D$~any point whatever
on the axis of abscissas, $AD = p$, $B$~any point whatever on the shortest line normal
to~$AD$ at~$D$, and $BD = q$, so that $p$~can be regarded as the abscissa, $q$~the ordinate
of the point~$B$. The abscissas we assume positive on the branch of the axis of
abscissas to which $\phi = 0$ corresponds, while we always regard~$r$ as positive. We take
the ordinates positive in the region in which $\phi$~is measured between $0$~and~$180°$.

By the theorem of \Art{16} we shall have
\[
\omega = 90°,\quad
F = 0,\quad
G = 1,
\]
and we shall set also
\[
\sqrt{E} = n.
\]
Thus $n$~will be a function of $p$,~$q$, such that for $q = 0$ it must become equal to unity.
The application of the formula of \Art{18} to our case shows that on any shortest
line \emph{whatever} we must have\Note{17}
\[
d\theta = \frac{\dd n}{\dd q}·dp,\NoteMark
\]
where $\theta$~denotes the angle between the element of this line and the element of the
line for which $q$~is constant. Now since the axis of abscissas is itself a shortest line,
and since, for it, we have everywhere $\theta = 0$, we see that for $q = 0$ we must have
everywhere
\[
\frac{\dd n}{\dd q} = 0.
\]
Therefore we conclude that, if $n$~is developed into a series in ascending powers of~$q$,
this series must have the following form:
\[
n = 1 + fq^{2} + gq^{3} + hq^{4} + \text{etc.}\Add{,}
\]
where $f$,~$g$,~$h$,~etc., will be functions of~$p$, and we set
\begin{alignat*}{4}
f &= f° &&+ f'p &&+ f''p^{2} &&+ \text{etc.}\Add{,} \\
g &= g° &&+ g'p &&+ g''p^{2} &&+ \text{etc.}\Add{,} \\
h &= h° &&+ h'p &&+ h''p^{2} &&+ \text{etc.}\Add{,}
\end{alignat*}
\PageSep{36}
or
\begin{alignat*}{2}
n = 1 + f°q^{2} &+ f'pq^{2}   &&+ f''p^{2}q^{2} + \text{etc.} \\
                  &+ g° q^{3} &&+ g'pq^{3} + \text{etc.} \\
                  &             &&+ h°q^{4} + \text{etc.\ etc.}
\end{alignat*}


\Article{24.}

The equations of \Art{22} give, in our case,
\begin{gather*}
 n\sin\psi = \frac{\dd r}{\dd p},\quad
  \cos\psi = \frac{\dd r}{\dd q},\quad
-n\cos\psi = m·\frac{\dd\phi}{\dd r},\quad
  \sin\psi = m·\frac{\dd\phi}{\dd q}, \\
%
n^{2} = n^{2}\left(\frac{\dd r}{\dd q}\right)^{2}
          + \left(\frac{\dd r}{\dd p}\right)^{2},\quad
n^{2}·\frac{\dd r}{\dd q}·\frac{\dd\phi}{\dd q}
      + \frac{\dd r}{\dd p}·\frac{\dd\phi}{\dd p} = 0\Add{.}
\end{gather*}
By the aid of these equations, the fifth and sixth of which are contained in the others,
series can be developed for $r$,~$\phi$,~$\psi$,~$m$, or for any functions whatever of these quantities.
We are going to establish here those series that are especially worthy of
attention.

Since for infinitely small values of $p$,~$q$ we must have
\[
r^{2} = p^{2} + q^{2},
\]
the series for~$r^{2}$ will begin with the terms $p^{2} + q^{2}$. We obtain the terms of higher
order by the method of undetermined coefficients,\footnote
  {We have thought it useless to give the calculation here, which can be somewhat abridged by
  certain artifices.}
by means of the equation
\[
\left(\frac{1}{n}·\frac{\dd(r^{2})}{\dd p}\right)^{2}
  + \left(\frac{\dd(r^{2})}{\dd q}\right)^{2} = 4r^{2}\Add{.}
\]
Thus we have\Note{18}
\begin{alignat*}{3}
\Tag{[1]}
r^{2} &= p^{2} + \tfrac{2}{3}f°p^{2}q^{2} &&+ \tfrac{1}{2}f'p^{3}q^{2} &&+ (\tfrac{2}{5}f'' - \tfrac{4}{45}{f°}^{2})p^{4}q^{2}\quad\text{etc.} \\
      &+ q^{2} &&+ \tfrac{1}{2}g°p^{2}q^{3} &&+ \tfrac{2}{5}g'p^{3}q^{3} \\
      &&& &&+(\tfrac{2}{5}h° - \tfrac{7}{45}{f°}^{2})p^{2}q^{4}\Add{.}
\end{alignat*}

Then we have, from the formula\Note{19}
\[
r\sin\psi = \frac{1}{2n}·\frac{\dd(r^{2})}{\dd p},
\]
\begin{alignat*}{2}
\Tag{[2]}
r\sin\psi = p - \tfrac{1}{3}f°pq^{2}
  & -\tfrac{1}{4}f'p^{2}q^{2} &&- (\tfrac{1}{5}f'' + \tfrac{8}{45}{f°}^{2})p^{3}q^{2}\quad\text{etc.} \\
  & -\tfrac{1}{2}g°pq^{3} &&-\tfrac{2}{5}g'p^{2}q^{3} \\
  &&& -(\tfrac{3}{5}h° - \tfrac{8}{45}{f°}^{2})pq^{4}\Add{;}
\end{alignat*}
\PageSep{37}
and from the formula\Note{20}
\[
r\cos\psi = \tfrac{1}{2}\, \frac{\dd(r^{2})}{\dd q}\Add{,}
\]
\begin{alignat*}{2}
\Tag{[3]}
r\cos\psi = q + \tfrac{2}{3}f°p^{2}q
  & +\tfrac{1}{2}f'p^{3}q &&+ (\tfrac{2}{5}f'' - \tfrac{4}{45}{f°}^{2})p^{4}q\quad\text{etc.} \\
  & +\tfrac{3}{4}g°p^{2}q^{2} &&+\tfrac{3}{5}g'p^{3}q^{2} \\
  &&& +(\tfrac{4}{5}h° - \tfrac{14}{45}{f°}^{2})p^{2}q^{3}\Add{.}
\end{alignat*}
These formulæ give the angle~$\psi$. In like manner, for the calculation of the angle~$\phi$,
series for $r\cos\phi$ and $r\sin\phi$ are very elegantly developed by means of the partial
differential equations
\begin{align*}
&\frac{\dd·r\cos\phi}{\dd p}
  = n\cos\phi·\sin\psi - r\sin\phi·\frac{\dd\phi}{\dd p}\Add{,} \\
&\frac{\dd·r\cos\phi}{\dd q}
  = \Z\cos\phi·\cos\psi - r\sin\phi·\frac{\dd\phi}{\dd q}\Add{,} \\
&\frac{\dd·r\sin\phi}{\dd p}
  = n\sin\phi·\sin\psi + r\cos\phi·\frac{\dd\phi}{\dd p}\Add{,} \\
&\frac{\dd·r\sin\phi}{\dd q}
  = \Z\sin\phi·\cos\psi + r\cos\phi·\frac{\dd\phi}{\dd q}\Add{,} \\
&n\cos\psi·\frac{\dd\phi}{\dd q} + \sin\psi·\frac{\dd\phi}{\dd p} = 0\Add{.}
\end{align*}
A combination of these equations gives
\begin{alignat*}{3}
&\frac{r\sin\psi}{n}·\frac{\dd·r\cos\phi}{\dd p}
  &&+ r\cos\psi·\frac{\dd·r\cos\phi}{\dd q}
  &&= r\cos\phi\Add{,} \\
&\frac{r\sin\psi}{n}·\frac{\dd·r\sin\phi}{\dd p}
  &&+ r\cos\psi·\frac{\dd·r\sin\phi}{\dd q}
  &&= r\sin\phi\Add{.}
\end{alignat*}
From these two equations series for $r\cos\phi$, $r\sin\phi$ are easily developed, whose first
terms must evidently be $p$,~$q$ respectively. The series are\Note{21}
\begin{alignat*}{3}
\Tag{[4]}
r\cos\phi &= p + \tfrac{2}{3}f°pq^{2}
  &&+ \tfrac{5}{12}f'p^{2}q^{2}
  &&+ (\tfrac{3}{10}f'' - \tfrac{8}{45}{f°}^{2})p^{3}q^{2}\quad\text{etc.} \\
 &&&+ \tfrac{1}{2}g°pq^{3}
  &&+ \tfrac{7}{20}g'p^{2}q^{3} \\
&&&&&+ (\tfrac{2}{5}h° - \tfrac{7}{45}{f°}^{2})pq^{4}\Add{,} \displaybreak[1] \\
\Tag{[5]}
r\sin\phi &= q - \tfrac{1}{3}f°p^{2}q
  &&- \tfrac{1}{6}f'p^{3}q
  &&- (\tfrac{1}{10}f'' - \tfrac{7}{90}{f°}^{2})p^{4}q\quad\text{etc.} \\
 &&&- \tfrac{1}{4}g°p^{2}q^{2}
  &&- \tfrac{3}{20}g'p^{3}q^{2} \\
&&&&&- (\tfrac{1}{5}h° + \tfrac{13}{90}{f°}^{2})p^{2}q^{3}\Add{.}
\end{alignat*}
From a combination of equations [2],~[3],~[4],~[5] a series for $r^{2}\cos(\psi + \phi)$, may
be derived, and from this, dividing by the series~[1], a series for $\cos(\psi + \phi)$, from
\PageSep{38}
which may be found a series for the angle $\psi + \phi$ itself. However, the same series
can be obtained more elegantly in the following manner. By differentiating the first
and second of the equations introduced at the beginning of this article, we obtain
\[
\sin\psi·\frac{\dd n}{\dd q}
  + n\cos\psi·\frac{\dd\psi}{\dd q}
  +  \sin\psi·\frac{\dd\psi}{\dd p} = 0\Add{,}
\]
and this combined with the equation
\[
n\cos\psi·\frac{\dd\phi}{\dd q} + \sin\psi·\frac{\dd\phi}{\dd p} = 0
\]
gives
\[
\frac{r\sin\psi}{n}·\frac{\dd n}{\dd q}
  + \frac{r\sin\psi}{n}·\frac{\dd(\psi + \phi)}{\dd p}
  + r\cos\psi·\frac{\dd(\psi + \phi)}{\dd q} = 0\Add{.}
\]
From this equation, by aid of the method of undetermined coefficients, we can easily
derive the series for $\psi + \phi$, if we observe that its first term must be~$\frac{1}{2}\pi$, the radius
being taken equal to unity and $2\pi$~denoting the circumference of the circle,\Note{22}
\begin{alignat*}{2}
\Tag{[6]}
\psi + \phi = \tfrac{1}{2}\pi - f°pq
  &- \tfrac{2}{3}f'p^{2}q
  &&- (\tfrac{1}{2}f'' - \tfrac{1}{6}{f°}^{2})p^{3}q\quad\text{etc.} \\
  &- g°pq^{2}
  &&- \tfrac{3}{4}g'p^{2}q^{2} \\
  &&&- (h° - \tfrac{1}{3}{f°}^{2})pq^{3}\Add{.}
\end{alignat*}

It seems worth while also to develop the area of the triangle~$ABD$ into a series.
For this development we may use the following conditional equation, which is easily
derived from sufficiently obvious geometric considerations, and in which $S$~denotes the
required area:\Note{23}
\[
\frac{r\sin\psi}{n}·\frac{\dd S}{\dd p}
  + r\cos\psi·\frac{\dd S}{\dd q}
  = \frac{r\sin\psi}{n}·\int n\, dq\Add{,}\NoteMark
\]
the integration beginning with $q = 0$. From this equation we obtain, by the method
of undetermined coefficients,\Note{24}
\begin{alignat*}{3}
\Tag{[7]}
S = \tfrac{1}{2}pq
  &- \tfrac{1}{12}f°p^{3}q
  &&- \tfrac{1}{20}f'p^{4}q
  &&- (\tfrac{1}{30}f'' - \tfrac{1}{60}{f°}^{2})p^{5}q\quad\text{etc.} \\
  &- \tfrac{1}{12}f°pq^{3}
  &&- \tfrac{3}{40}g°p^{3}q^{2}
  &&- \tfrac{1}{20}g'p^{4}q^{2} \\
  &&&- \tfrac{7}{120}f'p^{2}q^{3}
  &&- (\tfrac{1}{15}h° + \tfrac{2}{45}f'' + \tfrac{1}{60}{f°}^{2})p^{3}q^{3} \\
  &&&- \tfrac{1}{10}g°pq^{4}
  &&- \tfrac{3}{40}g'p^{2}q^{4} \\
  &&&&&- (\tfrac{1}{10}h° - \tfrac{1}{30}{f°}^{2})pq^{5}\Add{.}
\end{alignat*}
\PageSep{39}


\Article{25.}

From the formulæ of the preceding article, which refer to a right triangle formed
by shortest lines, we proceed to the general case. Let $C$ be another point on the
same shortest line~$DB$, for which point~$p$ remains the same as for the point~$B$, and
$q'$,~$r'$, $\phi'$,~$\psi'$, $S'$ have the same meanings as $q$,~$r$, $\phi$,~$\psi$, $S$ have for the point~$B$. There
will thus be a triangle between the points $A$,~$B$,~$C$, whose angles we denote by
$A$,~$B$,~$C$, the sides opposite these angles by $a$,~$b$,~$c$, and the area by~$\sigma$. We represent
the measure of curvature at the points $A$,~$B$,~$C$ by $\alpha$,~$\beta$,~$\gamma$ respectively. And then
supposing (which is permissible) that the quantities $p$,~$q$,~$q - q'$ are positive, we shall
have
\begin{align*}
A &= \phi - \phi', & B &= \psi, & C &= \pi - \psi', && \\
a &= q - q',       & b &= r',   & c &= r,           & \sigma &= S - S'.
\end{align*}

We shall first express the area~$\sigma$ by a series. By changing in~[7] each of the
quantities that refer to~$B$ into those that refer to~$C$, we obtain a formula for~$S'$.
Whence we have, exact to quantities of the sixth order,\Note{25}
\begin{align*}
\sigma = \tfrac{1}{2}p(q - q')
  \bigl(1 &- \tfrac{1}{6} f°(p^{2} + q^{2} + qq' + q'^{2}) \\
          &- \tfrac{1}{60}f'p(6p^{2} + 7q^{2} + 7qq' + 7q'^{2}) \\
          &- \tfrac{1}{20}g°(q + q')(3p^{2} + 4q^{2} + 4q'^{2})\bigr)\Add{.}\NoteMark
\end{align*}
This formula, by aid of series~[2], namely,
\[
c\sin B = p(1 - \tfrac{1}{3}f°q^{2}
              - \tfrac{1}{4}f'pq^{2}
              - \tfrac{1}{2}g°q^{3} - \text{etc.})
\]
can be changed into the following:
\begin{align*}
\sigma = \tfrac{1}{2}ac\sin B
  \bigl(1 &- \tfrac{1}{6} f°(p^{2} - q^{2} + qq' + q'^{2}) \\
          &- \tfrac{1}{60}f'p(6p^{2} - 8q^{2} + 7qq' + 7q'^{2}) \\
          &- \tfrac{1}{20}g°
    (3p^{2}q + 3p^{2}q' - 6p^{3} + 4q^{2}q' + 4qq'^{2} + 4q'^{3})\bigr)\Add{.}
\end{align*}

The measure of curvature for any point whatever of the surface becomes (by \Art{19},
where $m$,~$p$,~$q$ were what $n$,~$q$,~$p$ are here)
\begin{align*}
k &= -\frac{1}{n}·\frac{\dd^{2} n}{\dd q^{2}}
   = -\frac{2f + 6gq + 12hq^{2} + \text{etc.}}{1 + fq^{2} + \text{etc.}} \\
  &= -2f - 6gq - (12h - 2f^{2}) q^{2} - \text{etc.}
\end{align*}
Therefore we have, when $p$,~$q$ refer to the point~$B$,
\[
\beta = - 2f° - 2f'p - 6g°q - 2f''p^{2} - 6g'pq - (12h° - 2{f°}^{2})q^{2} - \text{etc.}
\]
\PageSep{40}
Also
\begin{align*}
\gamma &= -2f° - 2f'p - 6g°q' - 2f''p^{2} - 6g'pq' - (12h° - 2{f°}^{2})q'^{2} - \text{etc.}\Add{,} \\
\alpha &= -2f°\Add{.}
\end{align*}
Introducing these measures of curvature into the expression for~$\sigma$, we obtain the following
expression, exact to quantities of the sixth order (exclusive):\Note{26}
\begin{align*}
\sigma = \tfrac{1}{2}ac \sin B
  \bigl(1 &+ \tfrac{1}{120}\alpha(4p^{2} - 2q^{2} + 3qq' + 3q'^{2}) \\
          &+ \tfrac{1}{120}\beta (3p^{2} - 6q^{2} + 6qq' + 3q'^{2}) \\
          &+ \tfrac{1}{120}\gamma(3p^{2} - 2q^{2} + \Z qq' + 4q'^{2})\bigr)\Add{.}\NoteMark
\end{align*}
The same precision will remain, if for $p$,~$q$,~$q'$ we substitute $c\sin B$, $c\cos B$, $c\cos B - a$.
This gives\Note{27}
\begin{align*}
\Tag{[8]}
\sigma = \tfrac{1}{2}ac\sin B
  \bigl(1 &+ \tfrac{1}{120}\alpha(3a^{2} + 4c^{2} - \Z9ac\cos B) \\
          &+ \tfrac{1}{120}\beta (3a^{2} + 3c^{2} -  12ac\cos B) \\
          &+ \tfrac{1}{120}\gamma(4a^{2} + 3c^{2} - \Z9ac\cos B)\bigr)\Add{.}
\end{align*}
Since all expressions which refer to the line~$AD$ drawn normal to~$BC$ have disappeared
from this equation, we may permute among themselves the points $A$,~$B$,~$C$ and
the expressions that refer to them. Therefore we shall have, with the same precision,
\begin{align*}
\Tag{[9]}
\sigma = \tfrac{1}{2}bc\sin A
  \bigl(1 &+ \tfrac{1}{120}\alpha(3b^{2} + 3c^{2} -  12bc\cos A) \\
          &+ \tfrac{1}{120}\beta (3b^{2} + 4c^{2} - \Z9bc\cos A) \\
          &+ \tfrac{1}{120}\gamma(4b^{2} + 3c^{2} - \Z9bc\cos A)\bigr)\Add{,} \\
%
\Tag{[10]}
\sigma = \tfrac{1}{2}ab\sin C
  \bigl(1 &+ \tfrac{1}{120}\alpha(3a^{2} + 4b^{2} - \Z9ab\cos C) \\
          &+ \tfrac{1}{120}\beta (4a^{2} + 3b^{2} - \Z9ab\cos C) \\
          &+ \tfrac{1}{120}\gamma(3a^{2} + 3b^{2} -  12ab\cos C)\bigr)\Add{.}
\end{align*}


\Article{26.}

The consideration of the rectilinear triangle whose sides are equal to $a$,~$b$,~$c$ is of
great advantage. The angles of this triangle, which we shall denote by $A^{*}$,~$B^{*}$,~$C^{*}$,
differ from the angles of the triangle on the curved surface, namely, from $A$,~$B$,~$C$,
by quantities of the second order; and it will be worth while to develop these differences
accurately. However, it will be sufficient to show the first steps in these more
tedious than difficult calculations.

Replacing in formulæ [1],~[4],~[5] the quantities that refer to~$B$ by those that
refer to~$C$, we get formulæ for $r'^{2}$,~$r'\cos\phi'$, $r'\sin\phi'$. Then the development of the
expression
\PageSep{41}
\begin{align*}
  r^{2} + r'^{2} &- (q - q')^{2}
    - 2r\cos\phi·r'\cos\phi'
    - 2r\sin\phi·r'\sin\phi' \\
    &\quad= b^{2} + c^{2} - a^{2} - 2bc\cos A \\
    &\quad= 2bc(\cos A^{*} - \cos A),
\end{align*}
combined with the development of the expression
\[
r\sin\phi·r'\cos\phi' - r\cos\phi·r'\sin\phi' = bc\sin A,
\]
gives the following formula:
\begin{align*}
\cos A^{*} - \cos A
  = -(q - q')p\sin A
  \bigl(\tfrac{1}{3}f° &+ \tfrac{1}{6}f'p + \tfrac{1}{4}g°(q + q') \\
  &+ (\tfrac{1}{10}f'' - \tfrac{1}{45}{f°}^{2})p^{2}
   + \tfrac{3}{20}g'p(q + q') \\
  &+ (\tfrac{1}{5}h° - \tfrac{7}{90}{f°}^{2})(q^{2} + qq' + q'^{2})
   + \text{etc.}\bigr)
\end{align*}
From this we have, to quantities of the fifth order,\Note{28}
\begin{align*}
A^{*} - A = +(q - q')p
  \bigl(\tfrac{1}{3}f°
  &+ \tfrac{1}{6}f'p + \tfrac{1}{4}g°(q + q') + \tfrac{1}{10}f''p^{2} \\
  &+ \tfrac{3}{20}g'p(q + q') + \tfrac{1}{5}h°(q^{2} + qq' + q'^{2}) \\
  &- \tfrac{1}{90}{f°}^{2}(7p^{2} + 7q^{2} + 12qq' + 7q'^{2})\bigr)\Add{.}\NoteMark
\end{align*}
Combining this formula with
\[
2\sigma = ap\bigl(1 - \tfrac{1}{6}f°(p^{2} + q^{2} + qq' + q'^{2})- \text{etc.}\bigr)
\]
and with the values of the quantities $\alpha$,~$\beta$,~$\gamma$ found in the preceding article, we obtain,
to quantities of the fifth order,\Note{29}
\begin{align*}
\Tag{[11]}
A^{*} = A - \sigma\bigl(\tfrac{1}{6}\alpha
  &+ \tfrac{1}{12}\beta
   + \tfrac{1}{12}\gamma + \tfrac{2}{15}f''p^{2} + \tfrac{1}{5}g'p(q + q') \\
  &+ \tfrac{1}{5}h°(3q^{2} - 2qq' + 3q'^{2}) \\
  &+ \tfrac{1}{90}{f°}^{2}(4p^{2} - 11q^{2} + 14qq' - 11q'^{2})\bigr)\Add{.}
\end{align*}
By precisely similar operations we derive\Note{30}
\begin{align*}
\Tag{[12]}
B^{*} = B - \sigma\bigl(\tfrac{1}{12}\alpha
  &+ \tfrac{1}{6}\beta
   + \tfrac{1}{12}\gamma + \tfrac{1}{10}f''p^{2} + \tfrac{1}{10}g'p(2q + q') \\
  &+ \tfrac{1}{5}h°(4q^{2} - 4qq' + 3q'^{2}) \\
  &- \tfrac{1}{90}{f°}^{2}(2p^{2} + 8q^{2} - 8qq' + 11q'^{2})\bigr)\Add{,}
\displaybreak[1] \\
%
\Tag{[13]}
C^{*} = C - \sigma\bigl(\tfrac{1}{12}\alpha
  &+ \tfrac{1}{12}\beta
   + \tfrac{1}{6}\gamma + \tfrac{1}{10}f''p^{2} + \tfrac{1}{10}g'p(q + 2q') \\
  &+ \tfrac{1}{5}h°(3q^{2} - 4qq' + 4q'^{2}) \\
  &- \tfrac{1}{90}{f°}^{2}(2p^{2} + 11q^{2} - 8qq' + 8q'^{2})\bigr)\Add{.}
\end{align*}
From these formulæ we deduce, since the sum $A^{*} + B^{*} + C^{*}$ is equal to two right
angles, the excess of the sum~$A + B + C$ over two right angles, namely,\Note{31}
\begin{align*}
\Tag{[14]}
A + B + C= \pi + \sigma\bigl(\tfrac{1}{3}\alpha
  &+ \tfrac{1}{3}\beta
   + \tfrac{1}{3}\gamma + \tfrac{1}{3}f''p^{2} + \tfrac{1}{2}g'p(q + q') \\
  &+ (2h° - \tfrac{1}{3}{f°}^{2})(q^{2} - qq' + q'^{2})\bigr)\Add{.}
\end{align*}
This last equation could also have been derived from formula~[6].
\PageSep{42}


\Article{27.}

If the curved surface is a sphere of radius~$R$, we shall have
\[
\alpha = \beta = \gamma = -2f° = \frac{1}{R^{2}};\quad
f'' = 0,\quad
g' = 0,\quad
6h° - {f°}^{2} = 0,
\]
or
\[
h° = \frac{1}{24R^{4}}.
\]
Consequently, formula~[14] becomes
\[
A + B + C = \pi + \frac{\sigma}{R^{2}},
\]
which is absolutely exact. But formulæ [11],~[12],~[13] give
\begin{align*}
A^{*} &= A - \frac{\sigma}{3R^{2}}
  - \frac{\sigma}{180R^{4}}(2p^{2} - q^{2} + 4qq' - q'^{2})\Add{,} \\
B^{*} &= B - \frac{\sigma}{3R^{2}}
  + \frac{\sigma}{180R^{4}}(p^{2} - 2q^{2} + 2qq' + q'^{2})\Add{,} \\
C^{*} &= C - \frac{\sigma}{3R^{2}}
  + \frac{\sigma}{180R^{4}}(p^{2} + q^{2} + 2qq' - 2q'^{2})\Add{;}
\end{align*}
or, with equal exactness,
\begin{align*}
A^{*} &= A - \frac{\sigma}{3R^{2}}
  - \frac{\sigma}{180R^{4}}(b^{2} + c^{2} - 2a^{2})\Add{,} \\
B^{*} &= B - \frac{\sigma}{3R^{2}}
  - \frac{\sigma}{180R^{4}}(a^{2} + c^{2} - 2b^{2})\Add{,} \\
C^{*} &= C - \frac{\sigma}{3R^{2}}
  - \frac{\sigma}{180R^{4}}(a^{2} + b^{2} - 2c^{2})\Add{.}
\end{align*}
Neglecting quantities of the fourth order, we obtain from the above the well-known
theorem first established by the illustrious Legendre.


\Article{28.}

Our general formulæ, if we neglect terms of the fourth order, become extremely
simple, namely:
\begin{align*}
A^{*} &= A - \tfrac{1}{12}\sigma(2\alpha + \beta + \gamma)\Add{,} \\
B^{*} &= B - \tfrac{1}{12}\sigma(\alpha + 2\beta + \gamma)\Add{,} \\
C^{*} &= C - \tfrac{1}{12}\sigma(\alpha + \beta + 2\gamma)\Add{.}
\end{align*}
\PageSep{43}

Thus to the angles $A$,~$B$,~$C$ on a non-spherical surface, unequal reductions must
be applied, so that the sines of the changed angles become proportional to the sides
opposite. The inequality, generally speaking, will be of the third order; but if the
surface differs little from a sphere, the inequality will be of a higher order. Even in
the greatest triangles on the earth's surface, whose angles it is possible to measure,
the difference can always be regarded as insensible. Thus, \eg, in the greatest of
the triangles which we have measured in recent years, namely, that between the
points Hohehagen, Brocken, Inselberg, where the excess of the sum of the angles was
$14''. 85348$, the calculation gave the following reductions to be applied to the angles:
\setlength{\TmpLen}{2in}%
\begin{align*}
\Dotrow{Hohehagen}{-4''.95113\rlap{\Add{,}}} \\
\Dotrow{Brocken}{- 4''.95104\rlap{\Add{,}}} \\
\Dotrow{Inselberg}{-4''.95131\rlap{.}}
\end{align*}


\Article{29.}

We shall conclude this study by comparing the area of a triangle on a curved
surface with the area of the rectilinear triangle whose sides are $a$,~$b$,~$c$. We shall
denote the area of the latter by~$\sigma^{*}$; hence
\[
\sigma^{*} = \tfrac{1}{2}bc\sin A^{*}
           = \tfrac{1}{2}ac\sin B^{*}
           = \tfrac{1}{2}ab\sin C^{*}\Add{.}
\]

We have, to quantities of the fourth order,
\[
\sin A^{*} = \sin A - \tfrac{1}{12}\sigma\cos A·(2\alpha + \beta + \gamma)\Add{,}
\]
or, with equal exactness,
\[
\sin A = \sin A^{*}·\bigl(1 + \tfrac{1}{24}bc\cos A·(2\alpha + \beta + \gamma)\bigr)\Add{.}
\]
Substituting this value in formula~[9], we shall have, to quantities of the sixth order,
\begin{align*}
\sigma = \tfrac{1}{2}bc\sin A^{*}·\bigl(1
  &+ \tfrac{1}{120}\alpha(3b^{2}+ 3c^{2} - 2bc\cos A) \\
  &+ \tfrac{1}{120}\beta (3b^{2}+ 4c^{2} - 4bc\cos A) \\
  &+ \tfrac{1}{120}\gamma(4b^{2}+ 3c^{2} - 4bc\cos A)\bigr),
\end{align*}
or, with equal exactness,
\[
\sigma = \sigma^{*}\bigl(1
  + \tfrac{1}{120}\alpha(a^{2} + 2b^{2} + 2c^{2})
  + \tfrac{1}{120}\beta (2a^{2} + b^{2} + 2c^{2})
  + \tfrac{1}{120}\gamma(2a^{2} + 2b^{2} + c^{2})\Add{.}
\]
For the sphere this formula goes over into the following form:
\[
\sigma = \sigma^{*}\bigl(1 + \tfrac{1}{24}\alpha(a^{2} + b^{2} + c^{2})\bigr).
\]
\PageSep{44}
It is easily verified that, with the same precision, the following formula may be taken
instead of the above:
\[
\sigma
  = \sigma^{*}\sqrt{\frac{\sin A·\sin B·\sin C}
                         {\sin A^{*}·\sin B^{*}·\sin C^{*}}}\Add{.}
\]
If this formula is applied to triangles on non-spherical curved surfaces, the error, generally
speaking, will be of the fifth order, but will be insensible in all triangles such
as may be measured on the earth's surface.
\PageSep{45}


